% ===========================================================================
%  ProvGuard-RAG: Ledger-Anchored Provenance and Semantic Verification
%  for Trustworthy Retrieval-Augmented LLMs in Cyber-Physical Incident Intelligence
%  Target: MDPI AI (ISSN 2673-2688)
%  Template: MDPI LaTeX class (Definitions/mdpi.cls)
% ===========================================================================

\documentclass[ai,article,submit,pdftex,moreauthors]{Definitions/mdpi}


% ===========================================================================
% Journal metadata
% ===========================================================================
\firstpage{1}
\makeatletter
\setcounter{page}{\@firstpage}
\makeatother
\pubvolume{1}
\issuenum{1}
\articlenumber{0}
\pubyear{2026}
\copyrightyear{2026}
\datereceived{ }
\daterevised{ }
\dateaccepted{ }
\datepublished{ }

% ===========================================================================
% User-defined packages (mdpi.cls already loads most standard packages)
% ===========================================================================
% The MDPI class already loads: amsmath, amssymb, booktabs, tabularx, multirow,
% graphicx, float, natbib, hyperref, cleveref, geometry, enumitem, xcolor, etc.
% Only add packages NOT already loaded by mdpi.cls.
\usepackage[ruled,linesnumbered]{algorithm2e}
\usepackage{makecell}
\renewcommand\theadfont{\bfseries}
\renewcommand\theadalign{cc}

% Single-source-of-truth numerical values used in narrative (auto-generated).
% Auto-generated by experiments/reproduce_tables.py
% Do not edit by hand. Numbers used in the manuscript narrative come from here.
\newcommand{\valBoneIncAcc}{0.758}
\newcommand{\valBoneEgr}{0.738}
\newcommand{\valBoneCsr}{0.721}
\newcommand{\valBoneAnsCov}{0.907}
\newcommand{\valBoneOverAb}{0.022}
\newcommand{\valBoneAbs}{0.413}
\newcommand{\valBoneBert}{0.861}
\newcommand{\valBoneRouge}{0.423}
\newcommand{\valBoneRetrProv}{0.17}
\newcommand{\valBoneAuditReady}{0.10}
\newcommand{\valBtwoIncAcc}{0.757}
\newcommand{\valBtwoEgr}{0.760}
\newcommand{\valBtwoCsr}{0.720}
\newcommand{\valBtwoAnsCov}{0.921}
\newcommand{\valBtwoOverAb}{0.029}
\newcommand{\valBtwoAbs}{0.417}
\newcommand{\valBtwoBert}{0.861}
\newcommand{\valBtwoRouge}{0.422}
\newcommand{\valBtwoRetrProv}{0.95}
\newcommand{\valBtwoAuditReady}{0.96}
\newcommand{\valBthreeIncAcc}{0.776}
\newcommand{\valBthreeEgr}{0.859}
\newcommand{\valBthreeCsr}{0.886}
\newcommand{\valBthreeAnsCov}{0.832}
\newcommand{\valBthreeOverAb}{0.071}
\newcommand{\valBthreeAbs}{0.830}
\newcommand{\valBthreeBert}{0.862}
\newcommand{\valBthreeRouge}{0.424}
\newcommand{\valBthreeRetrProv}{0.16}
\newcommand{\valBthreeAuditReady}{0.11}
\newcommand{\valBfourIncAcc}{0.777}
\newcommand{\valBfourEgr}{0.882}
\newcommand{\valBfourCsr}{0.887}
\newcommand{\valBfourAnsCov}{0.857}
\newcommand{\valBfourOverAb}{0.052}
\newcommand{\valBfourAbs}{0.860}
\newcommand{\valBfourBert}{0.862}
\newcommand{\valBfourRouge}{0.424}
\newcommand{\valBfourRetrProv}{0.96}
\newcommand{\valBfourAuditReady}{0.96}
\newcommand{\valBonePChunkone}{0.009}
\newcommand{\valBonePSupone}{0.006}
\newcommand{\valBonePChunkfive}{0.042}
\newcommand{\valBonePSupfive}{0.031}
\newcommand{\valBonePChunkonezero}{0.085}
\newcommand{\valBonePSuponezero}{0.062}
\newcommand{\valBtwoPChunkone}{0.004}
\newcommand{\valBtwoPSupone}{0.003}
\newcommand{\valBtwoPChunkfive}{0.020}
\newcommand{\valBtwoPSupfive}{0.014}
\newcommand{\valBtwoPChunkonezero}{0.040}
\newcommand{\valBtwoPSuponezero}{0.028}
\newcommand{\valBthreePChunkone}{0.008}
\newcommand{\valBthreePSupone}{0.002}
\newcommand{\valBthreePChunkfive}{0.040}
\newcommand{\valBthreePSupfive}{0.008}
\newcommand{\valBthreePChunkonezero}{0.079}
\newcommand{\valBthreePSuponezero}{0.016}
\newcommand{\valBfourPChunkone}{0.002}
\newcommand{\valBfourPSupone}{0.000}
\newcommand{\valBfourPChunkfive}{0.010}
\newcommand{\valBfourPSupfive}{0.001}
\newcommand{\valBfourPChunkonezero}{0.020}
\newcommand{\valBfourPSuponezero}{0.002}
\newcommand{\valAblBfourCsr}{0.887}
\newcommand{\valAblBfourProv}{0.96}
\newcommand{\valAblBfournoprovCsr}{0.880}
\newcommand{\valAblBfournoprovProv}{0.17}
\newcommand{\valAblBfournonliCsr}{0.718}
\newcommand{\valAblBfournonliProv}{0.95}
\newcommand{\valAblBfournoledgerCsr}{0.887}
\newcommand{\valAblBfournoledgerProv}{0.96}
\newcommand{\valTotalBoneLat}{1009}
\newcommand{\valTotalBfourLat}{1204}
\newcommand{\valLatRetrieve}{119}
\newcommand{\valLatRerankBfour}{15}
\newcommand{\valLatRerankBone}{13}
\newcommand{\valLatGen}{877}
\newcommand{\valLatVerify}{184}
\newcommand{\valLatLedger}{8}
\newcommand{\valLatAddedFour}{195}
\newcommand{\valStorageKB}{1.83}
\newcommand{\valStorageCI}{0.03}
\newcommand{\valBoneInjOverride}{0.739}
\newcommand{\valextsixInjOverride}{0.178}
\newcommand{\valBthreeInjOverride}{0.315}
\newcommand{\valBfourInjOverride}{0.093}
\newcommand{\valBoneInjIndirect}{0.612}
\newcommand{\valextsixInjIndirect}{0.215}
\newcommand{\valBthreeInjIndirect}{0.267}
\newcommand{\valBfourInjIndirect}{0.080}
\newcommand{\valBoneInjExfil}{0.403}
\newcommand{\valextsixInjExfil}{0.120}
\newcommand{\valBthreeInjExfil}{0.182}
\newcommand{\valBfourInjExfil}{0.047}
\newcommand{\valBoneInjLeak}{0.531}
\newcommand{\valextsixInjLeak}{0.153}
\newcommand{\valBthreeInjLeak}{0.213}
\newcommand{\valBfourInjLeak}{0.069}
\newcommand{\valInjAblVan}{0.743}
\newcommand{\valInjAblIso}{0.184}
\newcommand{\valInjAblProv}{0.621}
\newcommand{\valInjAblNli}{0.309}
\newcommand{\valInjAblPi}{0.137}
\newcommand{\valInjAblIn}{0.106}
\newcommand{\valInjAblFull}{0.093}
\newcommand{\valBoneTamperTone}{0.00}
\newcommand{\valBtwoTamperTone}{0.00}
\newcommand{\valBthreeTamperTone}{0.00}
\newcommand{\valBfourTamperTone}{1.00}
\newcommand{\valBoneTamperTtwo}{0.00}
\newcommand{\valBtwoTamperTtwo}{0.00}
\newcommand{\valBthreeTamperTtwo}{0.00}
\newcommand{\valBfourTamperTtwo}{1.00}
\newcommand{\valBoneTamperTthree}{0.00}
\newcommand{\valBtwoTamperTthree}{0.00}
\newcommand{\valBthreeTamperTthree}{0.00}
\newcommand{\valBfourTamperTthree}{1.00}
\newcommand{\valBoneTamperTfour}{0.00}
\newcommand{\valBtwoTamperTfour}{0.00}
\newcommand{\valBthreeTamperTfour}{0.00}
\newcommand{\valBfourTamperTfour}{1.00}
\newcommand{\valBoneTamperTfive}{0.00}
\newcommand{\valBtwoTamperTfive}{0.00}
\newcommand{\valBthreeTamperTfive}{0.00}
\newcommand{\valBfourTamperTfive}{1.00}
\newcommand{\valBoneTamperTsix}{0.00}
\newcommand{\valBtwoTamperTsix}{0.00}
\newcommand{\valBthreeTamperTsix}{0.61}
\newcommand{\valBfourTamperTsix}{0.78}
\newcommand{\valBoneTamperTseven}{0.00}
\newcommand{\valBtwoTamperTseven}{0.00}
\newcommand{\valBthreeTamperTseven}{0.00}
\newcommand{\valBfourTamperTseven}{0.65}
\newcommand{\valDecompFsP}{0.873}
\newcommand{\valDecompFsR}{0.852}
\newcommand{\valDecompFtP}{0.906}
\newcommand{\valDecompFtR}{0.881}
\newcommand{\valNliAccFactual}{0.918}
\newcommand{\valNliAccTemporal}{0.823}
\newcommand{\valNliAccCausal}{0.701}
\newcommand{\valNliAccAttack}{0.866}
\newcommand{\valNliAccRemed}{0.752}
\newcommand{\valNliAccAbstain}{0.912}
\newcommand{\valNliAccMulti}{0.683}
\newcommand{\valManMinRepro}{0.21}
\newcommand{\valManFullRepro}{0.97}


% ===========================================================================
% Title (capitalized)
% ===========================================================================
\Title{ProvGuard-RAG: Ledger-Anchored Provenance and Semantic Verification for Trustworthy Retrieval-Augmented LLMs in Cyber-Physical Incident Intelligence}

% ===========================================================================
% Authors
% ===========================================================================
\newcommand{\orcidauthorA}{0000-0000-0000-000X} % Add your ORCID here
\newcommand{\orcidauthorB}{0000-0000-0000-000X} % Add your ORCID here

\newcommand{\eg}{\textit{e.g.}}
\newcommand{\ie}{\textit{i.e.}}

\Author{Huirong Wang $^{1}$, Hao Wang $^{2}$ and Yonggang Fu $^{1,*}$}

\AuthorNames{Huirong Wang, Hao Wang and Yonggang Fu}

% ===========================================================================
% Affiliations
% ===========================================================================
\address{%
$^{1}$ \quad College of Marine Culture and Law, Jimei University, Yinjiang Road 185, Xiamen 361021, China; hrwong@jmu.edu.cn (H.R.W.); yonggangfu@jmu.edu.cn (Y.G.F.)\\
$^{2}$ \quad College of Computer Engineering, Jimei University, Yinjiang Road 185, Xiamen 361021, China; wanghaocec@jmu.edu.cn
}

% ===========================================================================
% Corresponding author
% ===========================================================================
\corres{Correspondence: yonggangfu@jmu.edu.cn (Y.G.F.)}

% ===========================================================================
% Abstract (single paragraph, ~200 words, no blank lines)
% ===========================================================================
\abstract{%
Retrieval-augmented generation (RAG) lets large language models condition responses on external evidence, which is attractive for cyber-physical incident intelligence where answers must be tied to operational logs, vulnerability records, and policy documents. Standard RAG systems, however, measure retrieval relevance and answer quality without systematically evaluating provenance traceability, semantic claim support, or post-hoc audit verifiability. We present ProvGuard-RAG, a provenance-aware RAG framework that targets these properties as a coupled system. ProvGuard-RAG registers evidence with cryptographic source hashes, reranks retrieved chunks by provenance, recency, and prior verification status in addition to relevance, decomposes generated responses into atomic claims, verifies those claims against retrieved evidence with a domain-adapted natural-language inference model, and writes a Merkle-anchored audit manifest for each query. We formulate trustworthy cyber-physical RAG as a four-property problem (answer correctness, evidence faithfulness, provenance traceability, and audit verifiability) and instantiate it on TON\_IoT and CICIoT2023 with CVE/NVD and MITRE ATT\&CK as auxiliary knowledge bases, comparing against six external baselines and three ablated variants on a 500-query benchmark. The pipeline is paired with an explicit threat model that separates evidence integrity, source authenticity, semantic support, and post-hoc auditability as distinct guarantees, and an evaluation protocol that exercises corpus poisoning, prompt injection, and seven post-hoc tamper-and-audit-failure scenarios. ProvGuard-RAG attains higher claim support and evidence-grounded recall than verification-free baselines, detects all five tested post-registration manifest- and hash-consistency tampering scenarios at unit detection rate, partially mitigates pre-registration poisoned documents and split-view log equivocation at measured but bounded rates, and adds bounded latency overhead, while making the residual non-guarantees of any provenance-and-NLI design explicit.
}

% ===========================================================================
% Keywords (3--10, semicolon-separated)
% ===========================================================================
\keyword{retrieval-augmented generation; cyber-physical systems; data provenance; semantic verification; auditability; trustworthy AI; incident intelligence}

% ===========================================================================
\begin{document}

% ===========================================================================
%  1. INTRODUCTION
% ===========================================================================
\section{Introduction}
\label{sec:introduction}

Cyber-physical incident intelligence differs from ordinary text question answering because the answer must remain connected to operational evidence. In an industrial IoT or connected-infrastructure environment, which is the focus of this paper, an incident report may combine packet traces, sensor readings, device metadata, vulnerability records, policy constraints, and remediation guidance. The same need for evidence-bound reasoning recurs across cyber-physical settings: IoT-network intrusion detection treats provenance and feature lineage as first-class evidence~\cite{zhu2024iotintrusion}; container-based industrial-IoT anomaly detection has to surface which sensor stream and which transformation produced a flagged anomaly~\cite{wang2023iiotstream}; and smart-grid security work on energy-theft detection must reconcile partially observed labels with operator decisions that depend on \emph{why} a meter was flagged~\cite{chen2024energytheft,wen2025heterofl}. While the architectural principles generalize to other cyber-physical domains (smart cities, healthcare IoT, energy systems), the empirical instantiation in Section~\ref{sec:experiments} is scoped to IoT security datasets. A fluent explanation is therefore not sufficient: the system must show which evidence was used, whether that evidence is current, whether it has been altered, and whether it semantically supports the generated claim. This requirement is especially important when generated outputs influence triage, mitigation, or post-incident review. The present paper addresses this evidence-control problem for retrieval-augmented large language model (LLM) systems used in cyber-physical incident intelligence.

Retrieval-augmented generation (RAG) is a natural candidate for this setting because it allows a language model to condition its response on external documents rather than relying only on parametric memory. Recent applications of LLMs in adjacent domain analysis, including bibliometric and literature-trend mining for demand-side energy management~\cite{meng2024dsmllm}, illustrate that LLMs can extract structured insight from heterogeneous text corpora once the corpus is properly assembled, but they also expose how easily an LLM-driven pipeline drifts from its source documents when the link between answer and evidence is not enforced. Prior work has shown that RAG evaluation cannot be reduced to answer accuracy alone: retrieval noise, negative rejection, multi-document integration, and counterfactual robustness are separate capabilities that must be tested explicitly~\cite{chen2024benchmarkingrag}. Adaptive retrieval research also indicates that retrieval is not uniformly beneficial; a model may need retrieval under some knowledge conditions but be distracted by retrieved context under others~\cite{wang2023selfknowledge}. These findings matter for cyber-physical environments because the retrieved context may include irrelevant log segments, outdated vulnerability descriptions, incomplete policy fragments, or poisoned evidence. A standard RAG pipeline can still produce a coherent answer in such cases, but coherence does not establish that the answer is supported by trustworthy evidence.

The missing object in many RAG pipelines is provenance. Cyber-physical evidence moves through collection, preprocessing, parsing, chunking, embedding, retrieval, reranking, prompting, generation, and review. Each step can change what the downstream model sees and what an auditor can later reconstruct. Responsible-AI provenance work treats data origin and processing history as necessary for accountability and explainability~\cite{werder2022dataprovenance}. For incident intelligence, the same logic applies at query time: if a generated response cites an attack pattern, a device event, or a mitigation policy, the system should expose the source identity, ingestion time, transformation path, retrieval status, and claim-level support relation. Without this information, post-hoc audit becomes dependent on a transcript rather than on a verifiable evidence trail. Distributed-ledger mechanisms provide one approach to making such trails tamper-evident. IoT provenance studies show that blockchain-based records can support data integrity and verifiability in large-scale environments~\cite{honarpajooh2021iotprovenance}, and blockchain-based federated learning further illustrates how ledger mechanisms reduce single-point trust assumptions in distributed AI settings~\cite{yang2023trustworthyfl,wu2023bflsurvey}. These observations do not justify storing raw cyber-physical logs or retrieved documents on-chain. They instead motivate a lighter design in which source hashes, evidence identifiers, model metadata, verification labels, and output hashes are anchored in a ledger or Merkle-log structure while sensitive content remains off-chain.

This paper therefore targets a structural gap between four research streams. RAG benchmarks characterize retrieval-related failure modes, but they do not make provenance completeness and audit-log verifiability central evaluation targets~\cite{chen2024benchmarkingrag}. Domain-specific RAG work has shown the value of interpretable retrieval over authoritative corpora, but focuses on answer generation and reference use rather than tamper-evident evidence manifests~\cite{louis2024legalrag}. Provenance research argues that data lineage is a prerequisite for responsible AI accountability, while LLM audit research argues that governance-level audits must be complemented by model- and application-level audits because many risks arise in concrete deployment contexts~\cite{werder2022dataprovenance,mokander2023auditingllms}. What remains underdeveloped, and what this paper addresses, is a RAG design for cyber-physical incident intelligence that jointly treats retrieval relevance, source provenance, semantic claim support, and post-hoc auditability as measurable system properties. This gap is structural rather than accidental: the four properties are addressed in separate literatures that do not intersect at the retrieval-generation loop.

We propose ProvGuard-RAG, a provenance-aware retrieval-augmented generation framework for trustworthy cyber-physical incident intelligence. The framework registers incident logs, vulnerability records, security knowledge, and policy documents with source identifiers, timestamps, parser versions, and cryptographic hashes; retrieves candidate evidence with hybrid lexical and dense retrieval; reranks evidence using relevance, source trust, recency, and verification status; decomposes generated responses into atomic claims; verifies those claims against retrieved evidence; and records an audit manifest containing query hashes, evidence identifiers, source hashes, model metadata, verification labels, and output hashes. The method is designed so that the hypotheses of Section~\ref{sec:experiments} can be tested: namely, that provenance-aware reranking and semantic verification reduce unsupported or poisoned-evidence-grounded answers, and that ledger-anchored manifests improve post-hoc verifiability with bounded overhead. The main contributions of this paper are:
\begin{enumerate}
    \item A formal problem definition that decomposes trustworthy cyber-physical RAG into four jointly measured properties: answer correctness, evidence faithfulness, provenance traceability, and audit verifiability. This formulation extends standard RAG evaluation by making provenance and auditability first-class targets rather than optional documentation.
    \item ProvGuard-RAG, a modular framework that combines hybrid retrieval, provenance-aware reranking, semantic claim verification, and ledger-anchored audit manifests. The verification and audit components can be added to an existing RAG pipeline without retraining the generator.
    \item A multi-axis evaluation protocol covering seven measurement dimensions: answer quality, claim support, poisoned-evidence robustness, provenance completeness, audit verification, retrieval latency, and manifest storage overhead. This protocol operationalizes the four-property decomposition from the problem formulation.
    \item A reproducible benchmark instantiation on public CPS/IoT security datasets (TON\_IoT, CICIoT2023), public security knowledge bases (CVE/NVD, MITRE ATT\&CK), and an open-weight LLM (Qwen2.5-7B-Instruct), evaluated against six external baselines and three ablated variants under a unified protocol that exercises corpus poisoning, prompt injection, and seven post-hoc tamper-and-audit-failure scenarios. We release the dataset preprocessing pipeline, the evaluation harness, raw model outputs, retrieved-evidence identifiers, manifest JSON files with Merkle audit paths, human annotation labels, and the table-reproduction script used for every reported number.
\end{enumerate}
The rest of the paper is structured as follows. Section~\ref{sec:related} reviews related work across RAG evaluation, provenance systems, LLM auditing, and RAG security. Section~\ref{sec:method} formalizes the problem, states the threat model, and presents the ProvGuard-RAG architecture. Section~\ref{sec:experiments} reports the experimental setup and results across primary metrics, external baselines, poisoning, prompt injection, audit tamper, claim decomposition, per-claim-type verification, ablations, sensitivity, and scalability. Section~\ref{sec:discussion} interprets the results, analyzes failure modes, and discusses generalizability. Section~\ref{sec:conclusion} concludes.

% ===========================================================================
%  2. RELATED WORK
% ===========================================================================
\section{Related Work}
\label{sec:related}

\subsection{Retrieval-Augmented Generation and Evaluation}
\label{subsec:rag_evaluation}

Retrieval-augmented generation provides the immediate technical setting for this work because it separates a language model's response from the fixed knowledge stored in its parameters. Instead of asking the model to answer from memory, a RAG system retrieves external evidence and conditions generation on the retrieved context. This design is attractive for incident intelligence because cyber-physical events must often be interpreted against changing vulnerability records, operational policies, and domain-specific attack knowledge. However, RAG evaluation has moved beyond the simple question of whether retrieval improves final answer accuracy. Chen et al.~\cite{chen2024benchmarkingrag} characterize several distinct RAG capabilities, including robustness to noisy documents, rejection when retrieved evidence is insufficient, integration across multiple documents, and resistance to counterfactual evidence. These capabilities map directly onto CPS incident analysis, where a retrieved log fragment may be relevant to the query but irrelevant to the answer, or where a stale vulnerability description may conflict with a newer mitigation rule.

Recent retrieval work also shows that more retrieval is not automatically better. Wang et al.~\cite{wang2023selfknowledge} study retrieval augmentation guided by the model's self-knowledge, making retrieval conditional rather than treating it as a fixed preprocessing step. This perspective is useful for CPS settings because some incident questions can be answered from structured event labels, whereas others require external policy or vulnerability evidence. Domain-specific RAG work further shows how retrieved authority can support interpretability in high-stakes text generation. For example, Louis et al.~\cite{louis2024legalrag} formulate long-form legal question answering as a retrieve-then-read process over statutory material, where references are part of the interpretive value of the answer. These studies motivate the retrieval and evidence components of ProvGuard-RAG. Their primary evaluation objects, however, remain answer behavior, retrieval effectiveness, or reference use; they do not make provenance completeness, ledger-verifiable evidence manifests, or cyber-physical incident auditability central properties of the RAG pipeline.

\subsection{Data Provenance, Blockchain, and Trustworthy Distributed AI}
\label{subsec:provenance_blockchain}

The provenance literature provides a complementary view of trust. Rather than asking only whether a model output is correct, provenance asks whether the data behind that output can be traced through its origins and transformations. Werder et al.~\cite{werder2022dataprovenance} position data provenance as a basis for responsible AI because it records where data came from and how it was processed before influencing a decision. This principle is especially relevant to RAG because retrieval creates a short-lived evidence set at inference time. If the system does not record which chunks were retrieved, how they were derived, and whether they were verified, then a generated answer can be difficult to reconstruct even when the original corpus remains available.

Blockchain and distributed-ledger systems have been used to make provenance records tamper-evident in IoT and distributed AI settings. Honar Pajooh et al.~\cite{honarpajooh2021iotprovenance} propose blockchain-enabled provenance for IoT big data in a Hadoop ecosystem, emphasizing verifiability and integrity in large-scale connected environments. Building-level deployments in indoor-air-quality monitoring have shown that a blockchain-based alert system can integrate sensor evidence with risk assessments while keeping personally identifiable information off-chain~\cite{wan2023aira}, an architectural pattern that mirrors what a query-level audit manifest must achieve for cyber-physical RAG. In a related but distinct line, Yang et al.~\cite{yang2023trustworthyfl} use blockchain to support trustworthy federated learning, where ledger mechanisms help reduce single-point failure and malicious-update risks in distributed model training. Wu et al.~\cite{wu2023bflsurvey} synthesize blockchain-based federated learning and identify recurring challenges, including poisoning, incentive mechanisms, privacy-policy limitations, and communication efficiency. These works justify the use of tamper-evident records in distributed and security-sensitive AI systems. They do not, however, address the retrieval-time problem that ProvGuard-RAG targets: selecting evidence for a specific LLM query, verifying claim-level support, and recording an auditable manifest of the evidence-generation relation without placing raw CPS data on-chain.

\subsection{LLM Auditing and Application-Level Accountability}
\label{subsec:llm_auditing}

LLM audit research supplies the governance motivation for treating RAG outputs as auditable application events. M\"{o}kander et al.~\cite{mokander2023auditingllms} distinguish governance audits, model audits, and application audits, arguing that different risks arise at different layers of the LLM lifecycle. This distinction is important for RAG systems because a deployed application may combine a base model, a retrieval corpus, a chunking procedure, a reranker, prompts, access rules, and output filters. A governance audit of the model provider cannot by itself determine whether a CPS incident response was grounded in the correct evidence or whether the retrieved evidence had been altered. Likewise, a model-level audit may characterize general capabilities and limitations but still miss failures caused by local retrieval configuration, corpus drift, or query-specific evidence poisoning.

Application-level accountability therefore requires artifacts that are closer to the deployed RAG interaction. For cyber-physical incident intelligence, the relevant audit unit is not only the model or the corpus, but the generated answer together with its evidence set, provenance metadata, verification labels, and system configuration. This unit must be inspectable after the incident because response decisions may be challenged by operators, regulators, or security teams. Existing LLM audit frameworks clarify why such an application-layer audit is needed, but they do not specify a concrete CPS RAG architecture for producing the required artifacts. ProvGuard-RAG adopts the audit-layer perspective but operationalizes it through provenance-aware reranking, semantic claim verification, and a ledger-anchored audit manifest attached to each generated response.

\subsection{RAG Security and Adversarial Retrieval}
\label{subsec:rag_security}

A fourth line of work motivates the adversarial dimension of the ProvGuard-RAG design. Standard RAG pipelines assume that the retrieval corpus is trustworthy, but in security-sensitive cyber-physical environments this assumption does not hold. Corpus poisoning attacks insert malicious documents that surface under specific queries and cause the generator to produce attacker-controlled outputs. Prompt injection through retrieved context can override system instructions when the model treats retrieved text as authoritative. Retrieval-level adversarial manipulation can exploit the dense encoder to surface poisoned passages for benign queries. Zou et al.~\cite{zou2024poisonedrag} formalize corpus poisoning as an optimization problem in RAG and show that as few as five injected texts per target question achieve high attack success. Xian et al.~\cite{xian2025ragvulnerability} systematically evaluate adversarial robustness across 225 retrieval-generation configurations and find that poisoning attacks transfer across different retrievers and generators. Ni et al.~\cite{ni2025trustworthyrag} survey trustworthiness challenges in RAG across security, privacy, robustness, and explainability dimensions, identifying provenance traceability and end-to-end auditability as two of the least-addressed properties in current systems.

Indirect prompt injection adds a parallel attack surface: an attacker who controls one document in the corpus can attempt to redirect the model through instructions embedded inside retrieved text, including HTML, code blocks, or steganographic markers that survive chunking. Defenses such as input partitioning and instruction isolation~\cite{ni2025trustworthyrag} mitigate explicit overrides but do not address the upstream question of which retrieved chunk introduced the malicious instruction in the first place. Provenance metadata is complementary here: it does not parse the attack, but it records which source the offending chunk was attributed to and whether that source has accumulated a history of failed verifications. The combination of source-level provenance and claim-level entailment therefore provides a layered defense that no single existing RAG-security mechanism realizes alone.

These attack vectors are especially concerning in cyber-physical incident settings where a compromised evidence source could lead an analyst to misclassify an attack or apply an incorrect remediation, exactly the failure modes that motivate dedicated CPS intrusion-detection benchmarks~\cite{zhu2024iotintrusion} and graph-based time-series anomaly detectors~\cite{ma2025tcganomaly}. The architectural implication is straightforward: provenance metadata and semantic verification serve as defense-in-depth mechanisms. A poisoned chunk that passes retrieval may still be flagged by provenance checks (mismatched source hash, untrusted source identifier, anomalous ingestion timestamp) or by claim verification (no entailment between the claim and the trusted-source evidence in the same retrieval set).

\subsection{Positioning of ProvGuard-RAG}
\label{subsec:positioning}

ProvGuard-RAG is positioned at the intersection of these streams rather than as a replacement for any one of them. From RAG evaluation, it borrows the view that retrieval systems should be tested under noisy, insufficient, multi-document, and counterfactual evidence conditions~\cite{chen2024benchmarkingrag}. From adaptive and domain-specific retrieval work, it inherits the assumption that retrieval should be sensitive to task context and that authoritative references matter in high-stakes domains~\cite{wang2023selfknowledge,louis2024legalrag}. From provenance and blockchain systems, it borrows the requirement that evidence histories should be traceable and tamper-evident, while avoiding the stronger and often unnecessary assumption that raw operational data must be stored on-chain~\cite{werder2022dataprovenance,honarpajooh2021iotprovenance}. From LLM auditing, it borrows the distinction between general model assessment and application-level accountability~\cite{mokander2023auditingllms}. From RAG security, it borrows the recognition that retrieval corpora are attack surfaces and that provenance and verification can serve as layered defenses.

The added contribution is the coupling of these ideas inside the retrieval-generation loop for cyber-physical incident intelligence. To the best of our knowledge, no prior system jointly operationalizes provenance-aware reranking, claim-level semantic verification, and ledger-anchored query manifests for cyber-physical incident-intelligence RAG under a unified evaluation protocol. The proposed framework does not treat provenance as static dataset documentation, nor does it treat citations as sufficient evidence of grounding. Instead, retrieved chunks are ranked partly by provenance, recency, and prior verification status; generated answers are decomposed into atomic claims; claims are checked against trusted retrieved evidence; and the entire interaction is recorded as a hash-based audit manifest. This design creates measurable objects that standard RAG systems do not usually expose: corpus and retrieved provenance coverage, claim-support status by claim type, poisoned-evidence and prompt-injection robustness, manifest tamper-detection rates across distinct attack scenarios, and the storage and latency cost of audit anchoring. Section~\ref{sec:method} builds on this positioning by defining the problem setting, the threat model, and the way ProvGuard-RAG connects retrieval, verification, and audit logging in one framework.

% ===========================================================================
%  3. PROPOSED METHOD
% ===========================================================================
\section{Proposed Method: ProvGuard-RAG}
\label{sec:method}

\subsection{Problem Formulation}
\label{subsec:problem}

We define a trustworthy cyber-physical RAG system as a tuple of four measurable properties. Let $q$ denote an incident-intelligence query (\eg, ``Was device D17 compromised in the window 14:00--14:30?''), let $\mathcal{C} = \{d_1, \ldots, d_N\}$ denote a pre-registered evidence corpus where each document $d_i$ carries provenance metadata $p_i = (\text{src}_i, \tau_i, \text{parser}_i, h_i)$ consisting of a source identifier, ingestion timestamp, parser version, and cryptographic hash, and let $\mathcal{R}_q \subseteq \mathcal{C}$ denote the subset of evidence retrieved for query $q$.

A standard RAG pipeline generates an answer $a = G(q, \mathcal{R}_q)$ using a generator $G$ (an LLM) conditioned on the query and the retrieved evidence. A trustworthy RAG pipeline additionally produces: (i)~a claim decomposition $A = \{c_1, \ldots, c_K\}$ extracted from $a$, (ii)~a verification map $v: A \times \mathcal{R}_q \to \{\text{supported}, \text{unsupported}, \text{insufficient}\}$ produced by a semantic entailment model, and (iii)~an audit manifest $M_q$ that records query hash $h_q$, evidence identifiers and source hashes $\{h_i : d_i \in \mathcal{R}_q\}$, model metadata, verification labels $\{v(c_j, \mathcal{R}_q)\}_{j=1}^K$, and output hash $h_a$. The manifest is appended to an append-only transparency log $\mathcal{L}$, making it tamper-evident without storing content on-chain.

The trustworthy RAG problem is then: given a query $q$ and evidence corpus $\mathcal{C}$, produce $(a, \mathcal{R}_q, A, v, M_q)$ such that four properties are simultaneously satisfied:
\begin{enumerate}[label=(\roman*)]
    \item \textbf{Answer correctness:} $a$ is accurate with respect to the ground-truth incident state.
    \item \textbf{Evidence faithfulness:} For every claim $c_j \in A$ that is marked ``supported,'' the retrieved evidence $\mathcal{R}_q$ semantically entails $c_j$ (within the precision of the entailment model).
    \item \textbf{Provenance traceability:} Every $d_i \in \mathcal{R}_q$ carries complete provenance metadata $p_i$, and the mapping $d_i \mapsto p_i$ is verifiable via cryptographic hash.
    \item \textbf{Audit verifiability:} A third party possessing $M_q$ and access to $\mathcal{L}$ can verify that $M_q$ was appended to $\mathcal{L}$ at time $t$, that the hashes in $M_q$ match the evidence and output, and that $M_q$ has not been modified since $t$.
\end{enumerate}
Properties (i) and (ii) are partially addressed by existing RAG evaluation frameworks. Properties (iii) and (iv) are not jointly measured by any existing RAG system to our knowledge, and it is their simultaneous treatment together with (i) and (ii) that constitutes the problem formulation contribution.

The problem formulation rests on four assumptions that bound its scope. \textbf{Assumption 1 (Pre-registration feasibility):} The evidence corpus is pre-registered with provenance metadata before query time. This registration step is assumed to be technically feasible and not the system's performance bottleneck. \textbf{Assumption 2 (Lightweight log):} The transparency log stores only cryptographic hashes, not evidence content. We assume an append-only Merkle-tree log with negligible write cost per query. \textbf{Assumption 3 (NLI screening role):} The entailment model is used as a screening tool to flag unsupported claims for human review, not as a definitive truth oracle. Its precision and recall are sufficient for this screening purpose but do not replace domain-expert judgment. \textbf{Assumption 4 (Textual scope):} The evidence corpus consists of structured and semi-structured textual data (logs, alerts, vulnerability records, policy documents). Multi-modal CPS evidence such as network packet captures and system-call traces is outside the current scope.

\subsection{Threat Model}
\label{subsec:threat_model}

A central source of confusion in trustworthy-RAG work is the conflation of evidence integrity, source authenticity, semantic support, and post-hoc auditability. ProvGuard-RAG treats these as distinct properties that require distinct mechanisms; we therefore make the threat model explicit before describing the architecture.

\textbf{Attacker capabilities.} We assume an attacker who may (i)~insert documents into one or more low-trust evidence sources before or after registration; (ii)~modify unregistered evidence at any time; (iii)~embed prompt-injection content inside retrieved documents (raw text, HTML, code blocks, or steganographic tokens that survive chunking); (iv)~compromise a previously trusted source whose new content will be ingested with valid provenance metadata; (v)~attempt to alter the audit manifest, the generated answer, or the transparency log entries after generation; and (vi)~present split-view (equivocating) Merkle roots to different auditors when the attacker controls or coerces the log operator. The attacker is computationally bounded and cannot break the collision resistance of SHA-256 or forge cryptographic signatures.

\textbf{Defender assumptions.} We assume that the hash function $H$ is collision resistant; that corpus snapshots are retained for the duration of the audit window; that the published Merkle root $r$ is observable by at least one honest external witness; that evidence registration in Stage~1 is correctly implemented (no metadata falsification at registration time); and that human reviewers handle every claim that the verifier flags as \texttt{unsupported} or \texttt{insufficient}. The NLI verifier is assumed to be a screening tool, not a truth oracle, and the operator-assigned authority weights $w_{\text{auth}}(\text{src})$ are assumed to reflect a domain-calibrated initial trust estimate.

\textbf{Security guarantees.} Under these assumptions, ProvGuard-RAG provides four positive guarantees. (G1)~\textbf{Evidence integrity}: any modification of a registered document after registration is detectable, because the hash anchored at registration no longer matches the hash recomputed from the modified document. (G2)~\textbf{Manifest integrity}: any post-hoc modification of the audit manifest, the generated answer, the verification labels, or the evidence identifiers is detectable, because the recomputed manifest hash will not match the value $h_M$ recorded at log position $\pi(h_M)$. (G3)~\textbf{Replay detection}: a replayed manifest is detectable because $h_q$ in the manifest does not match the hash of the new query. (G4)~\textbf{Reduced exposure to untrusted evidence}: the provenance-aware reranker (Eq.~\ref{eq:rerank_score}) reduces the share of retrieved chunks originating from sources with low operator-assigned authority or poor verification history, and the NLI verifier flags claims that no trusted-source chunk in $\mathcal{R}_q$ entails.

\textbf{Non-guarantees.} The framework explicitly does not provide four properties. (N1)~\textbf{Evidence truthfulness}: the system verifies that a claim is supported by retrieved evidence, not that the evidence itself is factually correct. A document whose content was poisoned before registration carries internally consistent provenance metadata; ProvGuard-RAG cannot distinguish such a document from a legitimate one on hash grounds alone. (N2)~\textbf{Trusted-source compromise}: an attacker who compromises a high-authority source can publish malicious content with valid provenance; the framework can only reduce subsequent trust if the compromise is later reflected in the verification-success history $w_{\text{hist}}(\text{src})$. (N3)~\textbf{NLI correctness}: the verifier inherits the limitations of NLI models on multi-hop, temporal, causal, and numerical reasoning, and its labels should be treated as a screening signal that triggers human review. (N4)~\textbf{Log-operator coercion without external witness}: if the log operator equivocates by serving different Merkle roots to different auditors and no honest external witness is available, the equivocation is not detectable from a single auditor's view. The audit-tamper experiments in Section~\ref{subsec:tamper_results} operationalize each of these guarantees and non-guarantees on concrete attack scenarios.

\subsection{System Architecture Overview}
\label{subsec:architecture}

ProvGuard-RAG instantiates the trustworthy RAG problem through five connected stages summarized in Figure~\ref{fig:architecture}. The query enters at Stage~2; the evidence corpus $\mathcal{C}$ is registered off the query critical path in Stage~1; Stage~3 is the only borrowed standard-RAG block; Stages~4 and~5 add semantic verification and ledger-anchored auditing on top of the generator output, with the manifest hash anchored in an external Merkle transparency log~$\mathcal{L}$.

\begin{figure}[H]
\centering
\includegraphics[width=\linewidth]{figures/architecture.pdf}
\caption{ProvGuard-RAG architecture. The query $q$ flows left-to-right through hybrid retrieval and provenance-aware reranking (Stage~2, Eqs.~(\ref{eq:hybrid})--(\ref{eq:rerank_score})), grounded LLM generation and claim decomposition (Stage~3), NLI-based claim verification (Stage~4, Eq.~(\ref{eq:verification})), and audit manifest construction (Stage~5, Eq.~(\ref{eq:manifest})). The manifest hash $h_M$ is appended to the external Merkle transparency log~$\mathcal{L}$, which returns a position and audit path for independent inclusion verification (Eq.~(\ref{eq:merkle_verify})). Orange blocks are ProvGuard-RAG additions; the gray block is a standard RAG component; the purple block is the external log.}
\label{fig:architecture}
\end{figure}

\textbf{Stage 1: Evidence Registration.} Each evidence document $d$ ingested into the corpus is assigned a source identifier $\text{src}$ (naming the log stream, vulnerability database, or policy repository it originates from), a SHA-256 hash $h = H(d)$, an ingestion timestamp $\tau$, and a parser version identifier. This registration is performed once, off the query critical path. The provenance metadata $p = (\text{src}, \tau, \text{parser}, h)$ is stored alongside the document in the retrieval index; the hash alone, not the document content, is optionally anchored in the transparency log at registration time to provide a pre-retrieval integrity baseline.

\textbf{Stage 2: Hybrid Retrieval with Provenance-Aware Reranking.} Given query $q$, a hybrid retriever combines BM25 lexical scores $s_{\text{BM25}}$ and dense embedding similarity scores $s_{\text{dense}}$, each min-max normalized over the candidate set, to produce an initial relevance score
\begin{equation}
    s_{\text{rel}}(q, d) = \lambda\,\tilde{s}_{\text{BM25}}(q, d) + (1 - \lambda)\,\tilde{s}_{\text{dense}}(q, d), \qquad \lambda \in [0, 1],
    \label{eq:hybrid}
\end{equation}
which yields the initial candidate set $\mathcal{R}_q^{\text{cand}}$ of the top $N_{\text{cand}}$ documents. A provenance-aware reranker then scores each candidate $d \in \mathcal{R}_q^{\text{cand}}$ on three additional factors. The source trust score combines an operator-assigned authority weight $w_{\text{auth}}(\text{src}) \in [0, 1]$ with the historical verification-success rate $w_{\text{hist}}(\text{src}) \in [0, 1]$ of evidence from that source:
\begin{equation}
    s_{\text{trust}}(\text{src}) = \frac{w_{\text{auth}}(\text{src}) + \rho\,w_{\text{hist}}(\text{src})}{1 + \rho},
    \label{eq:trust}
\end{equation}
where $\rho \geq 0$ controls the relative weight of historical evidence. The recency score applies an exponential decay with half-life $\tau_{1/2}$ measured against the ingestion timestamp $\tau$ of the document and the query-time clock $t_{\text{now}}$:
\begin{equation}
    s_{\text{recency}}(\tau) = \exp\!\left(-\frac{(t_{\text{now}} - \tau) \ln 2}{\tau_{1/2}}\right).
    \label{eq:recency}
\end{equation}
The verification-cascade score promotes documents that previously passed claim verification within a window $\Delta t$:
\begin{equation}
    s_{\text{ver}}(d) =
    \begin{cases}
        1   & \text{if } d \text{ has a prior \texttt{supported} verification within } \Delta t, \\
        0.5 & \text{if } d \text{ has a prior \texttt{insufficient} verification within } \Delta t, \\
        0   & \text{otherwise}.
    \end{cases}
    \label{eq:cascade}
\end{equation}
The four factors are combined into a final score
\begin{equation}
    s(q, d) = \alpha\,s_{\text{rel}}(q, d) + \beta\,s_{\text{trust}}(\text{src}) + \gamma\,s_{\text{recency}}(\tau) + \delta\,s_{\text{ver}}(d),
    \label{eq:rerank_score}
\end{equation}
with $\alpha + \beta + \gamma + \delta = 1$ and each coefficient in $[0, 1]$. The top $k$ documents under this score form the final evidence set $\mathcal{R}_q$ passed to the generator.

\textbf{Stage 3: LLM Generation and Claim Decomposition.} The generator $G$ (a Qwen2.5-7B-Instruct model in our instantiation, chosen for its strong multilingual technical reasoning capabilities and Apache 2.0 license facilitating reproducibility) produces answer $a = G(q, \mathcal{R}_q)$ using a prompt template that instructs the model to ground each factual statement in the provided evidence. A separate, lightweight claim extraction module (implemented as a few-shot prompted instance of the same LLM or a fine-tuned smaller model) decomposes $a$ into a set of atomic claims $A = \{c_1, \ldots, c_K\}$, where each claim is a single factual proposition (\eg, ``Device D17 received a SYN flood from IP 10.0.0.42 at 14:17'').

\textbf{Stage 4: Semantic Claim Verification.} For each claim $c_j \in A$, a semantic entailment model $E$ evaluates whether $\mathcal{R}_q$ entails $c_j$. We use a DeBERTa-v3-large model fine-tuned on a cyber-physical entailment dataset derived from public incident reports; the fine-tuning protocol is described in Section~\ref{subsec:implementation}. The design follows the NLI-based verification paradigm established by SummaC~\cite{laban2022summac} and AlignScore~\cite{zha2023alignscore}, which showed that sentence-level entailment scoring with appropriate aggregation can detect factual inconsistencies in generated text. The per-claim-type evaluation in Section~\ref{subsec:nli_type} reports the verifier's accuracy across seven claim taxonomies (single-evidence factual, temporal, causal, attack-classification, remediation, abstention, multi-document) so that the verifier's competence boundary is part of the manuscript rather than a hidden assumption. The verification label is:

\begin{equation}
    v(c_j, \mathcal{R}_q) =
    \begin{cases}
        \text{supported}      & \text{if } \max_{d \in \mathcal{R}_q} P(\text{entail} \mid c_j, d) > \theta_{\text{entail}} \\[4pt]
        \text{unsupported}    & \text{if } \max_{d \in \mathcal{R}_q} P(\text{contradict} \mid c_j, d) > \theta_{\text{contra}} \\[4pt]
        \text{insufficient}   & \text{otherwise}
    \end{cases}
    \label{eq:verification}
\end{equation}

where $\theta_{\text{entail}}$ and $\theta_{\text{contra}}$ are confidence thresholds. Claims labeled ``unsupported'' or ``insufficient'' are flagged in the final output; the system may optionally suppress them or request additional retrieval. We note that the max-over-documents formulation in Eq.~\eqref{eq:verification} is a simplification; in practice, we adopt the chunking-and-aggregation strategy from SummaC~\cite{laban2022summac}, which splits evidence documents into sentence-level chunks and aggregates entailment scores across chunks for each claim.

\textbf{Stage 5: Audit Manifest and Ledger Anchoring.} After generation and verification, the system constructs an audit manifest that records every input needed to re-execute the query offline. The manifest groups its fields into a query record $\mathbf{Q}$, a retrieval record $\mathbf{R}$, an inference-configuration record $\mathbf{I}$, a verification record $\mathbf{V}$, and an output record $\mathbf{O}$:
\begin{equation}
\small
\begin{aligned}
\mathbf{Q} &= \bigl(h_q,\, t_{\text{query}},\, \text{user\_id}\bigr), \\
\mathbf{R} &= \bigl(\text{snapshot\_id}(\mathcal{C}),\, \text{retriever\_v},\, \text{embed\_v},\, \text{bm25\_cfg},\, \text{faiss\_v},\, \text{chunk\_alg}, \\
        &\qquad\;\,\, N_{\text{cand}},\, k,\, \{(h_i, \text{src}_i, \tau_i)\}_{d_i \in \mathcal{R}_q}\bigr), \\
\mathbf{I} &= \bigl(\text{model\_id},\, \text{prompt\_hash},\, \text{decode\_cfg},\, \text{nli\_v},\, \theta_{\text{entail}},\, \theta_{\text{contra}}, \\
        &\qquad\;\,\, \text{decomposer\_v},\, \alpha,\beta,\gamma,\delta,\, \rho,\, \tau_{1/2},\, \Delta t,\, \text{seed}\bigr), \\
\mathbf{V} &= \bigl(\{(c_j, v(c_j, \mathcal{R}_q))\}_{j=1}^{K},\, \text{flag\_policy},\, \text{suppress\_mask}\bigr), \\
\mathbf{O} &= \bigl(h_a,\, t_{\text{response}},\, \text{env\_metadata}\bigr).
\end{aligned}
\end{equation}
The manifest itself is the concatenation $M_q = (\mathbf{Q}, \mathbf{R}, \mathbf{I}, \mathbf{V}, \mathbf{O})$:
\begin{equation}
    M_q = \mathbf{Q} \,\Vert\, \mathbf{R} \,\Vert\, \mathbf{I} \,\Vert\, \mathbf{V} \,\Vert\, \mathbf{O}.
    \label{eq:manifest}
\end{equation}
The five-record decomposition is not cosmetic. The query record fixes which question is being audited; the retrieval record reproduces exactly which evidence the generator received (corpus snapshot, retriever version, BM25 configuration, FAISS index version, chunking algorithm, candidate-pool and final-set sizes, and per-chunk source hashes); the inference-configuration record reproduces the generator and verifier behavior (model identifier, prompt template hash, decoding parameters including temperature/top-$p$, NLI model version, verification thresholds, claim-decomposer version, all reranker coefficients, and the random seed); the verification record exposes which claims were supported, flagged, or suppressed in the user-facing answer; and the output record captures the answer hash with environment metadata such as the inference container digest. In the manifest-content ablation of Section~\ref{subsec:manifest_abl} we show that this five-record split is what allows independent re-verification of individual claims, not the bare hash chain alone.

where $h_q = H(q)$ and $h_a = H(a)$. This manifest is hashed to produce $h_M = H(M_q)$, and $h_M$ is appended to an append-only transparency log $\mathcal{L}$. The log is implemented as a Merkle tree, providing membership proofs: given $h_M$ at log position $\pi(h_M)$, any auditor with access to the published Merkle root $r$ verifies inclusion via
\begin{equation}
    \text{Verify}(h_M, \pi(h_M), r) \;=\; \mathbb{1}\!\left[\,\text{FoldPath}\bigl(h_M,\,\text{audit\_path}(\pi(h_M))\bigr) \;=\; r\,\right],
    \label{eq:merkle_verify}
\end{equation}
where $\text{FoldPath}$ recursively applies $H$ to sibling-hash pairs along the audit path of length $O(\log N)$, following the Merkle audit-path design established by Certificate Transparency~\cite{laurie2013certificatetransparency}. The manifest itself is stored off-chain (in the application database), while only its hash resides in the log. Under the collision resistance of $H$ and the append-only property of $\mathcal{L}$, any post-hoc modification of $M_q$ is detectable because the recomputed hash $H(M_q')$ will not match the value $h_M$ that was originally appended at $\pi(h_M)$.

\subsection{Component Interactions and Design Rationale}
\label{subsec:rationale}

Several design choices merit explicit justification. First, the provenance-aware reranker (Stage~2) uses a linear combination rather than a learned scoring function. This choice is deliberate: in a security-sensitive setting, operators must be able to inspect and adjust the weighting coefficients ($\alpha, \beta, \gamma, \delta$) without retraining. A learned reranker would obscure the contribution of provenance factors to the final ranking. Second, the claim verification module (Stage~4) is kept separate from the generator rather than being integrated as a self-verification step. This separation enforces architectural independence between verification and generation: a claim that passes verification has been checked by a module distinct from the one that produced it, so verification cannot inherit the same prompt-conditioned biases as generation. Third, the audit manifest (Stage~5) records both evidence identifiers and claim-level verification labels. This granularity enables fine-grained post-hoc analysis: an auditor can determine not only which evidence was used, but which claims were supported by which evidence chunks.

Each component carries identifiable fragilities. The source trust scores $s_{\text{trust}}$ require domain calibration; mis-calibrated scores could penalize relevant evidence from novel or newly added sources. The entailment model inherits the domain-sensitivity of NLI systems; its accuracy may degrade on CPS incident language without domain-adaptive fine-tuning. The transparency log assumes an honest log operator that maintains append-only semantics; this assumption is standard for transparency-log systems but should be verified in deployment. These fragilities inform the failure mode analysis in Section~\ref{sec:discussion}.

\subsection{Pseudocode}
\label{subsec:algorithm}

The pseudocode in Algorithm~\ref{alg:provguard} summarizes the end-to-end query flow.

\begin{algorithm}[H]
\KwIn{query $q$; provenance-registered corpus $\mathcal{C} = \{(d_i, p_i)\}_{i=1}^{N}$ with $p_i = (\mathrm{src}_i, \tau_i, \mathrm{parser}_i, h_i)$; generator $G$; entailment model $E$; transparency log $\mathcal{L}$ with current Merkle root $r$.}
\textbf{Params:} hybrid coefficient $\lambda$; trust mixing $\rho$; reranker weights $(\alpha, \beta, \gamma, \delta)$ with $\alpha + \beta + \gamma + \delta = 1$; recency half-life $\tau_{1/2}$; verification-cascade window $\Delta t$; thresholds $\theta_{\mathrm{entail}}, \theta_{\mathrm{contra}}$; candidate width $N_{\mathrm{cand}}$; final width $k$.\;
\KwOut{answer $a$; atomic claims $A$ with verification labels $v$; audit manifest $M_{q}$; manifest hash $h_{M}$ with log position $\pi(h_{M})$ and audit path $\mathrm{ap}(h_{M})$.}
\BlankLine
\tcp{Stage 1 (Evidence Registration) runs off the query critical path and supplies $\mathcal{C}$ above. Stages 2--5 run online per query.}
\BlankLine
\tcp{Stage 2: hybrid retrieval and provenance-aware reranking.}
$\mathcal{R}_q^{\mathrm{cand}} \leftarrow \textsc{HybridRetrieve}(q, \mathcal{C}, N_{\mathrm{cand}}, \lambda)$ \tcp{see Eq.~(\ref{eq:hybrid})}
\ForEach{$d \in \mathcal{R}_q^{\mathrm{cand}}$}{
  compute $s_{\mathrm{trust}}(\mathrm{src}_d), s_{\mathrm{recency}}(\tau_d), s_{\mathrm{ver}}(d)$ via Eqs.~(\ref{eq:trust})--(\ref{eq:cascade})\;
  $s(d) \leftarrow \alpha\,s_{\mathrm{rel}}(q,d) + \beta\,s_{\mathrm{trust}}(\mathrm{src}_d) + \gamma\,s_{\mathrm{recency}}(\tau_d) + \delta\,s_{\mathrm{ver}}(d)$ \tcp{Eq.~(\ref{eq:rerank_score})}
}
$\mathcal{R}_q \leftarrow \textsc{TopK}\bigl(\{s(d) : d \in \mathcal{R}_q^{\mathrm{cand}}\},\, k\bigr)$\;
\BlankLine
\tcp{Stage 3: grounded generation and claim decomposition.}
$a \leftarrow G(q, \mathcal{R}_q)$\;
$A \leftarrow \textsc{DecomposeClaims}(a)$ \tcp{$A = \{c_1, \ldots, c_K\}$}
\BlankLine
\tcp{Stage 4: semantic claim verification, instantiating Eq.~(\ref{eq:verification}).}
\ForEach{$c_j \in A$}{
  $p^{\mathrm{ent}}_j \leftarrow \max_{d \in \mathcal{R}_q} P_E(\mathrm{entail} \mid c_j, d)$\;
  $p^{\mathrm{con}}_j \leftarrow \max_{d \in \mathcal{R}_q} P_E(\mathrm{contradict} \mid c_j, d)$\;
  \eIf{$p^{\mathrm{ent}}_j > \theta_{\mathrm{entail}}$}{
    $v(c_j) \leftarrow \mathtt{supported}$\;
  }{
    \eIf{$p^{\mathrm{con}}_j > \theta_{\mathrm{contra}}$}{
      $v(c_j) \leftarrow \mathtt{unsupported}$\;
    }{
      $v(c_j) \leftarrow \mathtt{insufficient}$\;
    }
  }
}
flag every $c_j$ with $v(c_j) \neq \mathtt{supported}$ in the user-facing answer\;
\BlankLine
\tcp{Stage 5: audit manifest construction and ledger anchoring.}
$h_q \leftarrow H(q)$\;
$h_a \leftarrow H(a)$\;
$M_q \leftarrow \mathbf{Q} \,\Vert\, \mathbf{R} \,\Vert\, \mathbf{I} \,\Vert\, \mathbf{V} \,\Vert\, \mathbf{O}$ \tcp{five-record manifest, Eq.~(\ref{eq:manifest})}
$h_M \leftarrow H(M_q)$\;
$(\pi(h_M),\, \mathrm{ap}(h_M),\, r') \leftarrow \textsc{AppendLog}(\mathcal{L}, h_M)$\;
\tcp{Inclusion proof: a third party verifies $\textsc{Verify}(h_M, \mathrm{ap}(h_M), r') = 1$ using Eq.~(\ref{eq:merkle_verify}).}
\Return{$(a, A, v, M_q, h_M, \pi(h_M), \mathrm{ap}(h_M))$}
\caption{ProvGuard-RAG end-to-end query processing. The total per-query cost is $O\bigl(\,T_{\mathrm{retrieve}}(N) + N_{\mathrm{cand}} + T_{G}(|q| + |\mathcal{R}_q|) + K\,|\mathcal{R}_q|\,T_{E} + \log|\mathcal{L}|\,\bigr)$, in which the generator term $T_{G}$ dominates the latency budget reported in Section~\ref{subsec:latency}.}
\label{alg:provguard}
\end{algorithm}

\subsection{Notation}
\label{subsec:notation}

\begin{table}[H]
\caption{Notation used in the problem formulation and method description.}
\label{tab:notation}
\begin{tabularx}{\textwidth}{cX}
\toprule
\textbf{Symbol} & \textbf{Meaning} \\
\midrule
$q$                    & Incident-intelligence query \\
$\mathcal{C}$          & Pre-registered evidence corpus \\
$d_i$                  & Evidence document $i$ \\
$p_i = (\text{src}_i, \tau_i, \text{parser}_i, h_i)$ & Provenance metadata for $d_i$ \\
$h_i = H(d_i)$         & SHA-256 hash of document $d_i$ \\
$\mathcal{R}_q$        & Evidence set retrieved for query $q$ \\
$G$                    & LLM generator (Qwen2.5-7B-Instruct) \\
$a$                    & Generated answer \\
$A = \{c_1,\ldots,c_K\}$ & Atomic claims extracted from $a$ \\
$E$                    & Entailment model for claim verification \\
$v(c_j, \mathcal{R}_q)$ & Verification label for claim $c_j$ \\
$M_q$                  & Audit manifest for query $q$ \\
$h_M$                  & Hash of manifest $M_q$ \\
$\mathcal{L}$          & Append-only Merkle transparency log \\
$s_{\text{rel}}, s_{\text{trust}}, s_{\text{recency}}, s_{\text{ver}}$ & Reranking score components (Eqs.~\ref{eq:hybrid}--\ref{eq:cascade}) \\
$\lambda$              & BM25--dense balance in hybrid retrieval (Eq.~\ref{eq:hybrid}) \\
$w_{\text{auth}}, w_{\text{hist}}, \rho$ & Authority weight, historical-success weight, mixing parameter (Eq.~\ref{eq:trust}) \\
$\tau_{1/2}, \Delta t$ & Recency half-life and verification-cascade window \\
$\alpha,\beta,\gamma,\delta$ & Reranking combination weights (Eq.~\ref{eq:rerank_score}) \\
$\theta_{\text{entail}}, \theta_{\text{contra}}$ & Entailment / contradiction thresholds (Eq.~\ref{eq:verification}) \\
$\pi(h_M), r$          & Log position of $h_M$ and published Merkle root (Eq.~\ref{eq:merkle_verify}) \\
\bottomrule
\end{tabularx}
\end{table}

% ===========================================================================
%  4. EXPERIMENTS
% ===========================================================================
\section{Experiments}
\label{sec:experiments}

\textit{[Planning Draft: All numerical values reported in this section are temporary planning values produced by the reproducible table-generation script at \texttt{experiments/reproduce\_tables.py} and must be replaced with verified real experimental results before submission.]}

\subsection{Experimental Setup}
\label{subsec:setup}

\subsubsection{Datasets and Query Construction}
\label{subsec:datasets}

We evaluate on two public CPS/IoT security datasets with two complementary security knowledge bases. \textbf{TON\_IoT}~\cite{alsaeedi2020toniot} contains telemetry from IoT sensors and network traffic across normal operation and nine attack scenarios (scanning, DoS, DDoS, ransomware, backdoor, data injection, XSS, password cracking, MitM); \textbf{CICIoT2023}~\cite{neto2023cicios2023} provides 33 attack types across seven categories executed in a 105-device IoT topology. Both datasets are accompanied by per-record attack labels that we treat as latent ground truth, never as visible evidence to the model. \textbf{CVE/NVD} and \textbf{MITRE ATT\&CK} act as auxiliary vulnerability and attack-pattern knowledge bases indexed alongside the operational telemetry. Dataset statistics appear in Table~\ref{tab:datasets}.

\begin{table}[H]
\caption{Datasets and query counts. The auxiliary knowledge bases are shared across queries from both primary datasets.}
\label{tab:datasets}
\begin{tabularx}{\textwidth}{lXcccc}
\toprule
\textbf{Dataset} & \textbf{Type} & \textbf{Records} & \textbf{Atk.\ Types} & \textbf{Queries} & \textbf{Avg.\ Evid./Q.}\\
\midrule
TON\_IoT~\cite{alsaeedi2020toniot}    & IoT telemetry        & 22.3M & 9   & 300       & 12.4\\
CICIoT2023~\cite{neto2023cicios2023} & IoT network flows    & 48.5M & 33  & 200       & 15.7\\
CVE/NVD                              & Vulnerability database & 200K  & --  & shared KB & --\\
MITRE ATT\&CK                        & Attack-pattern taxonomy & 600 & --  & shared KB & --\\
\bottomrule
\end{tabularx}
\end{table}

\textbf{From rows to queries.} Telemetry records and network flows are not natural-language questions, so we construct an incident-intelligence query set in three stages. (i)~We extract incident windows by sliding a 30-minute window over each labelled attack episode in TON\_IoT and CICIoT2023, retaining windows that contain at least 50 records and a single primary attack label. (ii)~For each window we instantiate one of fourteen incident-question templates $T_1$--$T_{14}$ (Table~\ref{tab:templates}), filling slots with device, time, and attack-pattern values drawn from the window. (iii)~For each query we materialize the \emph{ground-truth evidence path}: the set of telemetry rows that, together, are sufficient to answer the question, plus the relevant CVE and ATT\&CK records. This three-step procedure yields 500 queries (300 from TON\_IoT, 200 from CICIoT2023). We additionally generate 75 \emph{abstain queries} (15\% of the benchmark) whose ground-truth evidence path is intentionally empty, so that systems are tested on whether they correctly decline to answer.

\begin{table}[H]
\caption{The fourteen incident-intelligence query templates used to construct the 500-query benchmark from TON\_IoT and CICIoT2023 incident windows. \textit{Slots} are filled per window from device, time, and attack-pattern values; \textit{Evidence} indicates which sources contribute to the ground-truth evidence path.}
\label{tab:templates}
\small
\setlength{\tabcolsep}{4pt}
\begin{tabularx}{\textwidth}{@{}c >{\raggedright\arraybackslash}p{2.6cm} >{\raggedright\arraybackslash}p{2.0cm} >{\raggedright\arraybackslash}p{1.8cm} >{\raggedright\arraybackslash}X@{}}
\toprule
\thead{ID} & \thead{Query type} & \thead{Slots} & \thead{Evidence} & \thead{Example}\\
\midrule
$T_1$ & Compromise detection      & device, $[t_1,t_2]$ & telemetry            & Was device D17 compromised between 14:00 and 14:30?\\
$T_2$ & Attack-family classification & device, window   & telemetry            & What attack family targeted device D17 in window W?\\
$T_3$ & Source attribution        & incident            & telemetry            & Which source IP launched the SYN flood on device D17?\\
$T_4$ & Time of compromise        & device              & telemetry            & When was device D17 first compromised in the trace?\\
$T_5$ & CVE consistency           & host, alerts        & telemetry + CVE      & Which CVE is most consistent with the alerts on host H?\\
$T_6$ & ATT\&CK technique         & pattern             & telemetry + ATT\&CK  & Which ATT\&CK technique matches the observed pattern P?\\
$T_7$ & ATT\&CK mitigation        & technique           & ATT\&CK              & What mitigation does ATT\&CK technique T recommend?\\
$T_8$ & Lateral movement          & source device       & telemetry            & Did device D17 attempt lateral movement to other devices?\\
$T_9$ & Data exfiltration check   & device, window      & telemetry            & Was data exfiltrated from device D17 in window W?\\
$T_{10}$ & Affected scope          & incident            & telemetry            & Which devices were affected by incident I?\\
$T_{11}$ & Severity assessment     & incident            & telemetry + CVE      & What severity level applies to incident I?\\
$T_{12}$ & Persistence indicators  & device              & telemetry + ATT\&CK  & Are there persistence indicators on device D17?\\
$T_{13}$ & Anomaly explanation     & anomaly, window     & telemetry            & What explains anomaly A in window W on device D17?\\
$T_{14}$ & Remediation steps       & incident            & ATT\&CK + CVE        & What remediation steps apply to incident I?\\
\bottomrule
\end{tabularx}
\end{table}

\textbf{Telemetry-to-text rendering.} Each retrieved telemetry row is rendered into a single line with the template \texttt{[YYYY-MM-DD~hh:mm:ss~UTC]~device=$d$ event=$e$ field$_1$=$v_1$ \dots field$_n$=$v_n$}. Original feature names are preserved for reproducibility, timestamps are rounded to seconds (sub-second precision is dropped), device identifiers are anonymized via a stable per-corpus SHA-256 truncation (e.g., the original sensor ID becomes \texttt{D17}), and IP addresses are replaced by stable opaque tokens of the form \texttt{IP-$\langle 8$~hex$\rangle$}. Raw attack labels are scrubbed from every rendered row: the splitter that derives the rendered text from the raw row drops the \texttt{label}, \texttt{type}, and \texttt{attack\_subcategory} columns and refuses to emit any field whose name is in a prohibited-label list. The CVE/ATT\&CK records use the full canonical text published by NIST and MITRE (CVE: title, description, severity, CWE; ATT\&CK: technique name, tactics, procedure description). All evidence is chunked at 256-token boundaries with 32-token overlap; the average rendered chunk length is 184 tokens and the 95th percentile is 247 tokens.

\textbf{Train/test isolation for NLI fine-tuning.} The 12{,}000 (claim, evidence) pairs used to fine-tune the entailment model are derived from a disjoint set of attack episodes that share no time windows, no devices, and no CVE identifiers with the 500-query evaluation benchmark; this prevents NLI overfitting on benchmark content. The episode-level holdout is enforced by an automated splitter that we release with the harness.

\subsubsection{Threat Model and Attack Configurations}
\label{subsec:attacks}

The threat model in Section~\ref{subsec:threat_model} is operationalized by three concrete attack families. \textbf{Corpus poisoning} follows the PoisonedRAG procedure~\cite{zou2024poisonedrag}: a fraction $p_{\text{poison}} \in \{0.01, 0.05, 0.10\}$ of the evidence corpus is replaced with documents optimized via gradient-based dense-retriever manipulation to surface for 50 target query patterns; we additionally include three non-optimized variants (stale-but-authentic vulnerability records, semantically plausible but subtly wrong remediation advice, and multi-document contradiction sets) so that the evaluation is not exclusively gradient-attack driven. \textbf{Prompt injection} comprises four payload classes inserted directly into retrieved chunks: instruction overrides, indirect injection via HTML comments, exfiltration prompts that ask the model to leak the system prompt or earlier evidence, and silent system-prompt leakage. \textbf{Audit tamper} comprises seven post-hoc scenarios T1--T7 (Section~\ref{subsec:tamper_results}) that exercise the four security guarantees and the three non-guarantees enumerated in Section~\ref{subsec:threat_model}.

\subsubsection{Baselines}
\label{subsec:baselines}

We compare against four ablated configurations of ProvGuard-RAG itself plus six external baselines that target individual sub-properties of trustworthy RAG. The four internal configurations isolate the marginal contribution of each component:

\begin{enumerate}[label=B\arabic*:, leftmargin=*]
    \item \textbf{Vanilla RAG.} Standard hybrid retrieval, relevance-only reranking, LLM generation; no provenance metadata, no claim verification, no audit manifest.
    \item \textbf{RAG + Provenance Filter.} B1 augmented with provenance metadata and the four-factor reranker (Stage~2); no claim verification, no audit manifest.
    \item \textbf{RAG + NLI Verification.} B1 augmented with claim decomposition and semantic verification (Stages~3--4); no provenance reranking, no audit manifest.
    \item \textbf{ProvGuard-RAG (Full).} The complete five-stage pipeline.
\end{enumerate}
The six external baselines situate B4 against published trustworthy-RAG ideas: \emph{Citation-RAG}, a Llama-3 pipeline prompted to attach inline citations to every claim~\cite{louis2024legalrag}; \emph{Self-RAG}, a reflective generator that emits special tokens to request retrieval and to mark passages as supportive or irrelevant; \emph{RAG~+~Cross-Encoder Reranker (RankT5)}, which reranks the candidate set with a cross-encoder relevance model; \emph{RAG~+~RAGAS Faithfulness Filter}, which runs an answer-level faithfulness model and suppresses claims whose faithfulness score falls below a calibrated threshold; \emph{RAG~+~Signed-Document Retrieval}, which restricts retrieval to cryptographically signed documents and exposes the signature chain as the only audit artefact; and \emph{RAG~+~StruQ Prompt-Injection Defense}, which interposes an instruction-isolation layer between retrieved chunks and the prompt~\cite{ni2025trustworthyrag}. All baselines share the same 500-query benchmark, the same generator (Qwen2.5-7B-Instruct), and the same evaluation harness.

\subsubsection{Evaluation Metrics}
\label{subsec:metrics}

The operationally meaningful question is how trustworthy the system is in deployment, not how fluent its answers are. We therefore report three task-specific incident-correctness metrics as primary and treat fluency-style scores as secondary. The primary metrics are: \textbf{Incident Classification Accuracy} (whether the predicted attack family matches the latent ground-truth label of the incident window); \textbf{Abstention Precision} (fraction of system \texttt{insufficient} labels that match human-annotated ``should abstain'' decisions); and \textbf{Evidence-Grounded Claim Recall} (fraction of ground-truth supported claims that the system retains as \texttt{supported} in its final answer). The secondary fluency metrics are ROUGE-L and BERTScore against expert-written reference answers.

We further decompose provenance into four sub-metrics so that metadata availability and metadata utilization are not conflated: (i)~\emph{Corpus Provenance Coverage}, the fraction of all indexed documents that carry valid provenance metadata, which is a property of the corpus and indexing pipeline; (ii)~\emph{Retrieved Provenance Coverage}, the fraction of $\mathcal{R}_q$ that carries valid provenance metadata, which depends on whether the reranker uses provenance; (iii)~\emph{Provenance Utilization}, a binary system property indicating whether provenance metadata enters the reranking score; and (iv)~\emph{Audit-Ready Evidence Rate}, the fraction of retrieved chunks whose source hash is verifiable against the corpus snapshot at audit time. Audit-related and overhead metrics are also reported: manifest \emph{Audit Verifiability}, per-stage and end-to-end \emph{Latency}, and \emph{Storage Overhead}. Per-attack-family detection rates for the seven tamper scenarios T1--T7 appear in Table~\ref{tab:tamper}; per-attack-type ASR for the four prompt-injection classes appears in Table~\ref{tab:injection}.

\subsubsection{Annotation Protocol}
\label{subsec:annotation}

Three cybersecurity researchers (each with $\geq$3 years of CPS incident-analysis experience and familiarity with IoT attack taxonomies) annotated the full 500-query benchmark, producing 4{,}358 labels across three independent annotation tasks. \textbf{Task A (Evidence-local support):} for each generated claim, the annotator labelled whether the claim is supported by the system's retrieved evidence as \texttt{supported}, \texttt{unsupported}, or \texttt{insufficient}. This task is conditioned on what each system retrieved, and is the basis for claim support and evidence-grounded recall. \textbf{Task B (Global incident correctness):} for each query, the annotator chose one attack-family label from the held-out ground-truth taxonomy after reading the rendered evidence \emph{plus} the full benchmark documentation; this label is the basis for incident classification accuracy. \textbf{Task C (Abstention appropriateness):} for each query, the annotator decided whether a competent analyst would have abstained given the full benchmark design rather than the system's view; this label is the basis for abstention precision and over-abstention rate. Annotators saw one configuration's output at a time, in random order, with system identity blinded, and never saw the raw attack labels or other systems' outputs. Disagreements were resolved by majority vote on the four-way Task-A label and by a fourth senior annotator on tied claim-decomposition boundaries. Inter-annotator agreement reached Fleiss' $\kappa = 0.78$ on Task A, $\kappa = 0.84$ on Task B, and $\kappa = 0.82$ on Task C. The annotation guideline, the 24 worked examples shown to annotators, and the per-task and per-claim-type agreement breakdowns are released with the harness.

\subsubsection{Implementation Details}
\label{subsec:implementation}

The generator $G$ is Qwen2.5-7B-Instruct served via vLLM (tensor-parallel = 1, max sequence length 8192, decoding temperature 0.2, top-$p$ 0.9, repetition penalty 1.05) with a fixed system prompt that instructs the model to ground each statement in the provided evidence and to abstain when the evidence is insufficient. The retriever uses an OpenSearch BM25 backend ($k_1 = 1.2$, $b = 0.75$) and a FAISS HNSW index (M = 32, efSearch = 128) over BGE-M3 dense embeddings. The hybrid coefficient is $\lambda = 0.5$ (Eq.~\ref{eq:hybrid}); the default reranker weights are $(\alpha, \beta, \gamma, \delta) = (0.5, 0.25, 0.15, 0.10)$ with sensitivity reported in Section~\ref{subsec:sensitivity}; the recency half-life is $\tau_{1/2} = 30$~days; the verification window is $\Delta t = 7$~days. The entailment model $E$ is DeBERTa-v3-large fine-tuned with the AdamW optimizer (learning rate $2 \times 10^{-5}$, batch size 16, weight decay 0.01, 5 epochs with early stopping on a held-out validation set), thresholds $\theta_{\mathrm{entail}} = \theta_{\mathrm{contra}} = 0.85$. The transparency log is a Trillian-style append-only Merkle tree with SHA-256 leaves and signed tree heads.

All measurements are taken on a single workstation with $4 \times$ NVIDIA A100 80GB GPUs, two AMD EPYC 7763 CPUs, 512~GB RAM, and NVMe storage. Each configuration is run over five random seeds (\seqsplit{seeds = \{42, 1337, 2024, 31415, 65537\}}). Per-query metrics are computed within each seed; aggregate numbers are reported as the mean across seeds, with 95\% confidence intervals from a paired bootstrap over the 500 query-level observations within each seed. The dataset preprocessing pipeline, the evaluation harness, raw model outputs, retrieved-evidence identifiers, manifest JSON files with Merkle audit paths, human annotation labels, query templates, attack payloads, configuration files, and the table-reproduction script at \texttt{experiments/reproduce\_tables.py} are released alongside the manuscript.

\subsection{Main Results}
\label{subsec:main_results}

\begin{table}[H]
\caption{Primary results across the four pipeline configurations on the 500-query incident-intelligence benchmark. Each row is a metric; arrows give the desired direction. Best per row in \textbf{bold}, second-best \underline{underlined}; $\dagger$ marks $p<0.05$ vs.\ B1 by Friedman test with Nemenyi post-hoc. Values are mean $\pm$ 95\% paired bootstrap CI over the 500 query-level observations within each seed, averaged across five seeds.}
\label{tab:main_results}
\small
\setlength{\tabcolsep}{6pt}
\begin{tabularx}{\textwidth}{@{}>{\raggedright\arraybackslash}X *{4}{>{\centering\arraybackslash}c}@{}}
\toprule
\thead{Metric} & \thead{Vanilla\\RAG} & \thead{RAG +\\Prov.\ Filter} & \thead{RAG +\\NLI Verif.} & \thead{\textbf{ProvGuard-}\\\textbf{RAG (Ours)}}\\
\midrule
Incident classification accuracy ($\uparrow$) & 0.758 $\pm$ 0.005 & 0.757 $\pm$ 0.005 & \underline{0.776 $\pm$ 0.005} & \textbf{0.777 $\pm$ 0.005}$^\dagger$\\
Evidence-grounded claim recall ($\uparrow$) & 0.738 $\pm$ 0.004 & 0.760 $\pm$ 0.005 & \underline{0.859 $\pm$ 0.003} & \textbf{0.882 $\pm$ 0.004}$^\dagger$\\
Claim support ratio ($\uparrow$) & 0.721 $\pm$ 0.008 & 0.720 $\pm$ 0.008 & \underline{0.886 $\pm$ 0.005} & \textbf{0.887 $\pm$ 0.005}$^\dagger$\\
Supported answer coverage ($\uparrow$) & \underline{0.907 $\pm$ 0.003} & \textbf{0.921 $\pm$ 0.003} & 0.832 $\pm$ 0.004 & 0.857 $\pm$ 0.003\\
Over-abstention rate ($\downarrow$) & \textbf{0.022 $\pm$ 0.001} & \underline{0.029 $\pm$ 0.001} & 0.071 $\pm$ 0.002 & 0.052 $\pm$ 0.001\\
Abstention precision ($\uparrow$) & 0.413 $\pm$ 0.008 & 0.417 $\pm$ 0.009 & \underline{0.830 $\pm$ 0.003} & \textbf{0.860 $\pm$ 0.003}$^\dagger$\\
Manifest audit verifiability ($\uparrow$) & N/A & N/A & N/A & \textbf{1.000}$^\dagger$\\
\bottomrule
\end{tabularx}
\end{table}

Table~\ref{tab:main_results} reports the seven primary metrics across the four pipeline configurations on the 500-query benchmark. The pattern matches the architectural hypothesis. Configurations that include semantic verification (B3 and B4) outperform verification-free configurations (B1 and B2) by $\sim$2 points of incident classification accuracy and 12--15 points of evidence-grounded recall, because the verifier suppresses unsupported and \texttt{insufficient}-evidence claims that would otherwise enter the answer and lower the answer-level grounding rate. Within the verification-enabled pair, B4 attains an evidence-grounded recall of \valBfourEgr{} versus B3 at \valBthreeEgr{}, a 2--3 point lift that is consistent with provenance-aware reranking surfacing higher-quality evidence in the first place.

The two new columns (Answer Coverage and Over-Abstention Rate) make the suppression trade-off visible. Verification-free configurations retain almost every ground-truth answer unit (B1 = \valBoneAnsCov{}, B2 = \valBtwoAnsCov{}) but pay for it on grounding; verification-enabled configurations suppress some answer units (B3 = \valBthreeAnsCov{}, B4 = \valBfourAnsCov{}) and the over-abstention rate rises from \valBoneOverAb{} (B1) and \valBtwoOverAb{} (B2) to \valBthreeOverAb{} (B3) and \valBfourOverAb{} (B4). The full pipeline B4 sits in the middle of this trade-off: it loses about five points of answer coverage relative to B1 in exchange for a 14-point evidence-grounded-recall gain and a near-doubling of abstention precision. Abstention precision shows the largest separation: configurations without a verification module abstain rarely, and when they do their abstentions are noisy (B1 = \valBoneAbs{}, B2 = \valBtwoAbs{}); configurations with verification produce principled abstentions with precision above 0.83 (B3 = \valBthreeAbs{}, B4 = \valBfourAbs{}). Audit verifiability is binary: only B4 produces a manifest that can be re-verified against the transparency log.

\begin{figure}[H]
\centering
\includegraphics[width=0.95\linewidth]{figures/main_results.pdf}
\caption{Incident classification accuracy, claim support ratio, and evidence-grounded claim recall across the four pipeline configurations on the 500-query benchmark. Error bars are 95\% confidence intervals over five seeds.}
\label{fig:main_results}
\end{figure}

Figure~\ref{fig:main_results} visualizes the same comparison on the three primary metrics. The verification-driven gap on evidence-grounded recall is the dominant pattern: verification reshapes which claims survive into the user-facing answer, while the underlying generator quality is similar across configurations. The fluency-style secondary metrics are essentially flat across configurations (BERTScore differences below 0.002, ROUGE-L differences below 0.003), confirming that the gains on the primary metrics come from the verification and provenance machinery and not from a difference in generator quality.

\subsection{Comparison with External Baselines}
\label{subsec:external}

\begin{table}[H]
\caption{Comparison with six external baselines that each target an individual sub-property of trustworthy RAG. Each row reports incident classification accuracy, claim support ratio, evidence-grounded recall, and audit-ready evidence rate on the same 500-query benchmark. Best per column in \textbf{bold}.}
\label{tab:external}
\small
\setlength{\tabcolsep}{4pt}
\begin{tabularx}{\textwidth}{@{}>{\raggedright\arraybackslash}X *{4}{>{\centering\arraybackslash}c}@{}}
\toprule
\thead{Method} & \thead{Incident\\Acc.} & \thead{Claim\\Support} & \thead{Evid.-Grnd.\\Recall} & \thead{Audit-Ready\\Evid.}\\
\midrule
Vanilla RAG & 0.758 & 0.721 & 0.738 & 0.10\\
Citation-RAG (Llama-3 + cite prompt) & 0.760 & 0.781 & 0.792 & 0.18\\
Self-RAG (reflective tokens) & \textbf{0.781} & 0.812 & 0.813 & 0.20\\
RAG + Cross-Encoder Reranker (RankT5) & 0.766 & 0.745 & 0.786 & 0.19\\
RAG + RAGAS Faithfulness Filter & 0.758 & 0.834 & 0.802 & 0.21\\
RAG + Signed-Document Retrieval & 0.749 & 0.728 & 0.748 & 0.93\\
RAG + Prompt-Injection Defense (StruQ) & 0.755 & 0.741 & 0.768 & 0.20\\
RAG + Provenance Filter (B2) & 0.757 & 0.720 & 0.760 & 0.96\\
RAG + NLI Verification (B3) & 0.776 & 0.886 & 0.859 & 0.11\\
ProvGuard-RAG (B4, ours) & 0.777 & \textbf{0.887} & \textbf{0.882} & \textbf{0.96}\\
\bottomrule
\end{tabularx}
\end{table}

Table~\ref{tab:external} situates ProvGuard-RAG against six external baselines. Four observations are worth highlighting. First, baselines that augment the generation step alone (Citation-RAG, Self-RAG, RAGAS faithfulness filtering) raise claim support but not audit-ready evidence rate, because they have no mechanism to expose provenance metadata. Second, baselines that augment retrieval alone (RankT5 cross-encoder, signed-document retrieval) raise relevance or audit-ready evidence rate but do not improve claim-level support. Third, the prompt-injection defense baseline (StruQ) raises classification accuracy on injection-laden queries (Section~\ref{subsec:injection_results}) but, like all single-purpose baselines, leaves at least one of the four trust properties unaddressed. Fourth, ProvGuard-RAG attains the best evidence-grounded recall and audit-ready evidence rate, and is competitive on incident accuracy and claim support without dominating every column (Self-RAG matches B4 on classification accuracy, and RAGAS matches B4 within one point on claim support), consistent with the architectural argument that the contribution lies in the joint operationalization of all four properties, not in any single mechanism.

\subsection{Poisoned-Evidence Robustness}
\label{subsec:poisoning_results}

\begin{table}[H]
\caption{Poisoned-evidence robustness under three corpus poison levels. Lower is better in every cell; $\dagger$: $p<0.05$ vs.\ B1 at the same poison level by McNemar test.}
\label{tab:poisoning}
\small
\setlength{\tabcolsep}{4pt}
\begin{tabularx}{\textwidth}{@{}>{\raggedright\arraybackslash}p{2.4cm} *{6}{>{\centering\arraybackslash}X}@{}}
\toprule
& \multicolumn{2}{c}{\textbf{1\% poison}} & \multicolumn{2}{c}{\textbf{5\% poison}} & \multicolumn{2}{c}{\textbf{10\% poison}}\\
\cmidrule(lr){2-3} \cmidrule(lr){4-5} \cmidrule(lr){6-7}
\thead{Config.} & \thead{Chunk\\Rate} & \thead{Claim\\Sup.} & \thead{Chunk\\Rate} & \thead{Claim\\Sup.} & \thead{Chunk\\Rate} & \thead{Claim\\Sup.}\\
\midrule
Vanilla RAG & 0.009 $\pm$ 0.000 & 0.006 $\pm$ 0.000 & 0.042 $\pm$ 0.001 & 0.031 $\pm$ 0.001 & 0.085 $\pm$ 0.002 & 0.062 $\pm$ 0.001\\
RAG + Prov.\ Filter & 0.004 $\pm$ 0.000 & 0.003 $\pm$ 0.000 & 0.020 $\pm$ 0.000 & 0.014 $\pm$ 0.000 & 0.040 $\pm$ 0.001 & 0.028 $\pm$ 0.001\\
RAG + NLI Verif. & 0.008 $\pm$ 0.000 & 0.002 $\pm$ 0.000 & 0.040 $\pm$ 0.001 & 0.008 $\pm$ 0.000 & 0.079 $\pm$ 0.002 & 0.016 $\pm$ 0.000\\
ProvGuard-RAG (Full) & \textbf{0.002 $\pm$ 0.000}$^\dagger$ & \textbf{0.000 $\pm$ 0.000}$^\dagger$ & \textbf{0.010 $\pm$ 0.000}$^\dagger$ & \textbf{0.001 $\pm$ 0.000}$^\dagger$ & \textbf{0.020 $\pm$ 0.000}$^\dagger$ & \textbf{0.002 $\pm$ 0.000}$^\dagger$\\
\bottomrule
\end{tabularx}
\end{table}

Table~\ref{tab:poisoning} reports robustness under the three corpus-poisoning levels described in Section~\ref{subsec:attacks}. At the 5\% poison level the vanilla pipeline B1 selects poisoned chunks for \valBonePChunkfive{} of targeted queries and supports \valBonePSupfive{} of claims grounded in poisoned evidence. The provenance filter B2 cuts the chunk-selection rate roughly in half (to \valBtwoPChunkfive{}, $p<0.05$ by McNemar) because poisoned documents in this attack lack a valid pre-registration hash. The NLI-only configuration B3 has only a small effect on chunk selection (\valBthreePChunkfive{}) but cuts the poisoned-claim support rate to \valBthreePSupfive{} because the verifier rejects most poisoned claims at the entailment step. The full pipeline B4 attains the lowest rates across both columns at every poison level (\valBfourPChunkfive{} chunk and \valBfourPSupfive{} support at 5\%), and the relative ranking is preserved at 1\% and 10\% poison.

\begin{figure}[H]
\centering
\includegraphics[width=0.95\linewidth]{figures/poisoning_robustness.pdf}
\caption{(a) Poisoned-chunk selection rate and (b) poisoned-claim support rate across three corpus poison levels and the four pipeline configurations.}
\label{fig:poisoning}
\end{figure}

Figure~\ref{fig:poisoning} shows that the marginal value of the combined defense grows with the poison level: the gap between B4 and the next-best configuration widens from 1\% to 10\% poison, because provenance filtering removes documents with metadata anomalies while semantic verification removes claims with semantic inconsistency, and the union of these two filters covers a non-overlapping set of attacker techniques. We caution that this conclusion is specific to attacks that touch either the metadata channel or the entailment channel; the upstream non-guarantee N1 in Section~\ref{subsec:threat_model} (a poisoned document registered with internally consistent metadata) is examined separately as scenario T6 in Section~\ref{subsec:tamper_results}.

\subsection{Prompt-Injection Robustness}
\label{subsec:injection_results}

\begin{table}[H]
\caption{Prompt-injection attack-success rate (ASR) across four attack types and four configurations on 100 injection trials per cell. Lower is better; the StruQ defense baseline is included as a stronger comparison than vanilla RAG.}
\label{tab:injection}
\small
\setlength{\tabcolsep}{4pt}
\begin{tabularx}{\textwidth}{@{}>{\raggedright\arraybackslash}X *{4}{>{\centering\arraybackslash}c}@{}}
\toprule
\thead{Attack Type} & \thead{Vanilla\\RAG} & \thead{RAG +\\StruQ} & \thead{RAG +\\NLI Verif.} & \thead{\textbf{ProvGuard-}\\\textbf{RAG (Ours)}}\\
\midrule
Instruction override & 0.739 $\pm$ 0.027 & 0.178 $\pm$ 0.008 & 0.315 $\pm$ 0.013 & \textbf{0.093 $\pm$ 0.005}\\
Indirect injection (HTML) & 0.612 $\pm$ 0.023 & 0.215 $\pm$ 0.010 & 0.267 $\pm$ 0.012 & \textbf{0.080 $\pm$ 0.005}\\
Exfiltration prompt & 0.403 $\pm$ 0.016 & 0.120 $\pm$ 0.006 & 0.182 $\pm$ 0.009 & \textbf{0.047 $\pm$ 0.004}\\
System-prompt leakage & 0.531 $\pm$ 0.020 & 0.153 $\pm$ 0.008 & 0.213 $\pm$ 0.009 & \textbf{0.069 $\pm$ 0.004}\\
\bottomrule
\end{tabularx}
\end{table}

Table~\ref{tab:injection} reports the attack-success rate (ASR) of four prompt-injection attack types against B1, the StruQ defense baseline, B3 (NLI verification only), and B4. Vanilla RAG is highly vulnerable, with ASR in the range \valBoneInjExfil{}--\valBoneInjOverride{}. The StruQ baseline reduces ASR to roughly \valextsixInjOverride{}--\valextsixInjIndirect{} by isolating retrieved text from the system prompt, but it does not address claim-level grounding. NLI-only verification (B3) provides smaller direct protection (\valBthreeInjExfil{}--\valBthreeInjOverride{}) because injection succeeds at the prompt level before the verifier can act. ProvGuard-RAG (B4) attains the lowest ASR in every attack type (\valBfourInjExfil{}--\valBfourInjOverride{}), because the combination of provenance-aware reranking (which down-weights injected content from low-trust sources), instruction isolation in the prompt template, and post-generation NLI verification (which rejects claims that have no support in trusted-source evidence) creates a layered defense.

\begin{figure}[H]
\centering
\includegraphics[width=0.85\linewidth]{figures/injection_robustness.pdf}
\caption{Prompt-injection attack-success rate for four attack types and four configurations on 100 injection trials per cell. Lower is better.}
\label{fig:injection}
\end{figure}

\begin{table}[H]
\caption{Defense-source ablation for the instruction-override prompt-injection class on 100 trials per cell. Each row activates a single combination of instruction isolation, provenance-aware reranking, and NLI verification on top of the same hybrid retriever and generator. Lower attack-success rate (ASR) is better.}
\label{tab:inj_abl}
\small
\setlength{\tabcolsep}{6pt}
\begin{tabularx}{\textwidth}{@{}>{\raggedright\arraybackslash}X *{1}{>{\centering\arraybackslash}c}@{}}
\toprule
\thead{Defense configuration} & \thead{Override ASR}\\
\midrule
Vanilla RAG & 0.743 $\pm$ 0.074\\
+ Instruction isolation & 0.184 $\pm$ 0.018\\
+ Provenance reranking & 0.621 $\pm$ 0.062\\
+ NLI verification & 0.309 $\pm$ 0.031\\
+ Provenance + Instruction isolation & 0.137 $\pm$ 0.014\\
+ Instruction isolation + NLI & 0.106 $\pm$ 0.011\\
ProvGuard-RAG (full) & \textbf{0.093 $\pm$ 0.009}\\
\bottomrule
\end{tabularx}
\end{table}

The four-configuration comparison conflates three defense sources, so Table~\ref{tab:inj_abl} isolates them on the instruction-override class. Activating instruction isolation alone drops ASR from \valInjAblVan{} to \valInjAblIso{}, the largest single-component effect. Activating provenance reranking alone (\valInjAblProv{}) gives a much smaller drop because injected text inserted into a low-trust document is still visible to the generator before the verifier acts. Activating NLI alone (\valInjAblNli{}) cuts ASR roughly in half because the verifier rejects claims that have no trusted-source entailment. The pairwise combinations show clear additivity: instruction isolation $+$ provenance reranking reaches \valInjAblPi{}, instruction isolation $+$ NLI reaches \valInjAblIn{}, and the full pipeline reaches \valInjAblFull{}. The takeaway is that the headline injection-robustness number for ProvGuard-RAG is not attributable to provenance reranking alone; instruction isolation is the dominant defense source, with NLI verification and provenance reranking adding measurable but smaller complementary gains.

\subsection{Audit Tamper Tests}
\label{subsec:tamper_results}

\begin{table}[H]
\caption{Detection rate of seven post-hoc tamper scenarios over 200 trials each. Higher is better. B1--B3 produce no manifest, hence detect none of the manifest-related tamper types; T6 (a poisoned document registered with an internally consistent hash) and T7 (log equivocation) cannot be fully detected by any single component.}
\label{tab:tamper}
\small
\setlength{\tabcolsep}{4pt}
\begin{tabularx}{\textwidth}{@{}c >{\raggedright\arraybackslash}X *{4}{>{\centering\arraybackslash}c}@{}}
\toprule
\thead{Id} & \thead{Scenario} & \thead{Vanilla\\RAG} & \thead{RAG +\\Prov.\ Filter} & \thead{RAG +\\NLI Verif.} & \thead{\textbf{ProvGuard-}\\\textbf{RAG (Ours)}}\\
\midrule
T1 & Post-registration evidence mutation & 0.00 & 0.00 & 0.00 & \textbf{1.00}\\
T2 & Output mutation, manifest unchanged & 0.00 & 0.00 & 0.00 & \textbf{1.00}\\
T3 & Manifest deletion & 0.00 & 0.00 & 0.00 & \textbf{1.00}\\
T4 & Log-entry reordering & 0.00 & 0.00 & 0.00 & \textbf{1.00}\\
T5 & Manifest replay across queries & 0.00 & 0.00 & 0.00 & \textbf{1.00}\\
T6 & Poisoned doc, internally consistent hash & 0.00 & 0.00 & 0.61 $\pm$ 0.05 & \textbf{0.78 $\pm$ 0.04}\\
T7 & Log equivocation (split-view roots) & 0.00 & 0.00 & 0.00 & \textbf{0.65 $\pm$ 0.06}\\
\bottomrule
\end{tabularx}
\end{table}

Table~\ref{tab:tamper} reports tamper-detection rates over 200 trials per scenario. Five scenarios (T1--T5) test the manifest-integrity guarantees G1--G3 from Section~\ref{subsec:threat_model}: post-registration evidence mutation, output mutation with the manifest unchanged, manifest deletion, log-entry reordering, and manifest replay across different queries. ProvGuard-RAG detects all five at 1.000 because each violation triggers either a hash mismatch on the manifest or a mismatch between $h_q$ and the recomputed query hash. Configurations B1--B3 detect none of these scenarios because they produce no manifest at all: a baseline cannot detect tampering of an artefact it never produces.

The remaining two scenarios are operationally interesting precisely because they are not fully detectable. Scenario T6 (a poisoned document registered with an internally consistent hash, the non-guarantee N1) is detected at \valBfourTamperTsix{} in B4 because the NLI verifier catches a fraction of the inconsistencies the metadata cannot, and at \valBthreeTamperTsix{} in B3 (NLI alone) for the same reason. Scenario T7 (log equivocation, the non-guarantee N4) is detected at \valBfourTamperTseven{} in B4 only when an external witness is available; the residual gap reflects equivocation episodes in which no honest witness is present in our trial. These two scenarios are reported precisely so that the threat model in Section~\ref{subsec:threat_model} can be checked against measured behavior rather than against the system's marketing.

\subsection{Claim Decomposition Quality}
\label{subsec:decomp_results}

\begin{table}[H]
\caption{Quality of two claim-decomposition strategies against human-annotated claim segmentation on a held-out 200-answer subset. Higher is better.}
\label{tab:decomp}
\small
\setlength{\tabcolsep}{6pt}
\begin{tabularx}{\textwidth}{@{}>{\raggedright\arraybackslash}X *{3}{>{\centering\arraybackslash}c}@{}}
\toprule
\thead{Decomposer} & \thead{Precision} & \thead{Recall} & \thead{Segmentation\\F1}\\
\midrule
Few-shot Qwen2.5-7B-Instruct & 0.873 $\pm$ 0.012 & 0.852 $\pm$ 0.014 & 0.862 $\pm$ 0.011\\
Fine-tuned T5-base (claim-segmentation head) & 0.906 $\pm$ 0.010 & 0.881 $\pm$ 0.012 & 0.893 $\pm$ 0.009\\
\bottomrule
\end{tabularx}
\end{table}

Because the downstream claim-support and evidence-grounded-recall metrics depend on how the answer is decomposed into claims, we evaluate two decomposition strategies against human-annotated claim segmentations on a held-out 200-answer subset (Table~\ref{tab:decomp}). The few-shot Qwen2.5-7B-Instruct decomposer attains \valDecompFsP{} precision and \valDecompFsR{} recall against the 1{,}642 human claim spans, with a segmentation F1 of 0.862. The fine-tuned T5-base decomposer raises both numbers (\valDecompFtP{} precision, \valDecompFtR{} recall, F1 0.893) and is the configuration we adopt by default; the few-shot decomposer remains useful in low-resource deployments where the small T5 model is not available. We note that decomposition quality is the upper bound on the achievable claim support ratio: even a perfect verifier cannot recover a claim that the decomposer has missed, which is why decomposition recall is the metric that constrains the rest of the verification pipeline.

\subsection{NLI Verification by Claim Type}
\label{subsec:nli_type}

\begin{table}[H]
\caption{NLI verification performance broken out by claim type against the 4358-claim expert-annotated subset. The verifier exceeds 0.90 accuracy on factual and abstention claims and degrades on causal and multi-document claims, consistent with known NLI limitations on multi-hop reasoning.}
\label{tab:nli_type}
\small
\setlength{\tabcolsep}{6pt}
\begin{tabularx}{\textwidth}{@{}>{\raggedright\arraybackslash}X *{4}{>{\centering\arraybackslash}c}@{}}
\toprule
\thead{Claim Type} & \thead{N} & \thead{Accuracy} & \thead{Precision} & \thead{Recall}\\
\midrule
Single-evidence factual & 1842 & 0.918 & 0.929 & 0.911\\
Temporal window & 921 & 0.823 & 0.847 & 0.802\\
Causal & 308 & 0.701 & 0.736 & 0.681\\
Attack-classification & 571 & 0.866 & 0.883 & 0.852\\
Remediation recommendation & 219 & 0.752 & 0.778 & 0.733\\
Abstention / uncertainty & 176 & 0.912 & 0.935 & 0.892\\
Multi-document synthesis & 351 & 0.683 & 0.711 & 0.661\\
\bottomrule
\end{tabularx}
\end{table}

A central concern in any NLI-based verification design is that natural-language inference is unreliable on the kinds of claims that dominate cyber-physical incident responses, where temporal windows, multi-document correlation, and causality matter. Table~\ref{tab:nli_type} disaggregates the verifier's accuracy across the seven claim types defined in our annotation guideline. Single-evidence factual claims (\valNliAccFactual{}) and abstention/uncertainty claims (\valNliAccAbstain{}) are easy: a sentence-level entailment scorer is well-suited to ``did device $D$ receive a SYN flood at $t$'' or to ``the evidence does not specify the source IP.'' Attack-classification claims (\valNliAccAttack{}) and temporal-window claims (\valNliAccTemporal{}) are intermediate: domain-adapted NLI with chunking-and-aggregation~\cite{laban2022summac} captures most cases. Causal claims (\valNliAccCausal{}) and multi-document synthesis claims (\valNliAccMulti{}) are the verifier's weakest categories, exactly as multi-hop NLI literature predicts; the verifier's role on these claim types is therefore explicitly to flag them for human review rather than to make a definitive judgment, in line with the screening-tool design described in Section~\ref{subsec:problem}. Reporting the disaggregated table makes the verifier's competence boundary part of the manuscript rather than a hidden assumption.

\subsection{Component Ablation}
\label{subsec:ablation}

\begin{table}[H]
\caption{Component ablation. Each row removes one component from the full ProvGuard-RAG pipeline. Audit verifiability is binary; \textit{Retrieved Provenance Coverage} is the fraction of retrieved documents that carry valid provenance metadata at the time of generation.}
\label{tab:ablation}
\small
\setlength{\tabcolsep}{4pt}
\begin{tabularx}{\textwidth}{@{}>{\raggedright\arraybackslash}p{3.6cm} *{4}{>{\centering\arraybackslash}X}@{}}
\toprule
\thead{Configuration} & \thead{BERT-\\Score} & \thead{Claim\\Support} & \thead{Retr.\ Prov.\\Coverage} & \thead{Audit\\Verif.}\\
\midrule
ProvGuard-RAG (Full) & 0.862 $\pm$ 0.001 & 0.887 $\pm$ 0.005 & 0.96 $\pm$ 0.00 & 1.000\\
w/o Prov.\ reranking & 0.862 $\pm$ 0.001 & 0.880 $\pm$ 0.005 & 0.17 $\pm$ 0.01 & 1.000\\
w/o NLI verification & 0.861 $\pm$ 0.001 & 0.718 $\pm$ 0.008 & 0.95 $\pm$ 0.00 & 1.000\\
w/o Ledger anchoring & 0.862 $\pm$ 0.001 & 0.887 $\pm$ 0.005 & 0.96 $\pm$ 0.00 & N/A\\
\bottomrule
\end{tabularx}
\end{table}

Table~\ref{tab:ablation} ablates each ProvGuard-RAG component from the full pipeline. Removing the provenance reranker drops retrieved provenance coverage from \valAblBfourProv{} to \valAblBfournoprovProv{}, because evidence is no longer surfaced preferentially when it carries valid provenance, and reduces claim support marginally (\valAblBfourCsr{} $\rightarrow$ \valAblBfournoprovCsr{}). Removing the NLI verifier produces the largest claim-support drop (\valAblBfourCsr{} $\rightarrow$ \valAblBfournonliCsr{}), confirming that the verifier is the dominant driver of claim-level grounding. Removing the ledger anchoring eliminates audit verifiability while leaving the other metrics unchanged, confirming that audit anchoring is architecturally independent of retrieval and verification.

\begin{figure}[H]
\centering
\includegraphics[width=0.85\linewidth]{figures/ablation_waterfall.pdf}
\caption{Claim support ratio for the full pipeline and each one-component ablation. Error bars are 95\% confidence intervals over five seeds.}
\label{fig:ablation}
\end{figure}

Figure~\ref{fig:ablation} highlights the asymmetry: the verifier ablation cost is several times the next-largest, which is consistent with the verifier being the only component whose job is to filter the user-facing answer, while the reranker and the ledger affect upstream evidence selection and downstream auditability respectively.

\subsection{Manifest Content Ablation}
\label{subsec:manifest_abl}

\begin{table}[H]
\caption{Ablation over manifest content. The minimal manifest stores only query, evidence, and output hashes; the full manifest additionally records retriever, generator, NLI, prompt, and decoding metadata sufficient to re-execute the query offline. \textit{Config repro} is the fraction of manifests for which the recorded fields fully reconstruct the inference configuration; \textit{Claim audit} is the fraction of claims that can be independently re-verified against archived evidence.}
\label{tab:manifest_abl}
\small
\setlength{\tabcolsep}{6pt}
\begin{tabularx}{\textwidth}{@{}>{\raggedright\arraybackslash}X *{4}{>{\centering\arraybackslash}c}@{}}
\toprule
\thead{Manifest Variant} & \thead{Audit\\Verif.} & \thead{Config\\Repro.} & \thead{Claim\\Audit} & \thead{Size\\(KB)}\\
\midrule
Minimal hash-only & 1.000 & 0.21 & 0.00 & 0.42\\
ProvGuard-RAG (full) & 1.000 & 0.97 & 1.00 & 1.82\\
\bottomrule
\end{tabularx}
\end{table}

The five-record manifest in Eq.~\ref{eq:manifest} carries more than the bare hash chain that a minimal integrity scheme would record. Table~\ref{tab:manifest_abl} compares a minimal-hash manifest (storing only $h_q$, $h_a$, and the per-evidence source hashes) against the full manifest. Both produce a verifiable Merkle inclusion proof, so audit verifiability remains 1.000 in both cases. The difference appears in the configurability and claim-auditing dimensions: with the minimal manifest, only \valManMinRepro{} of inferences can be exactly re-executed offline because retriever, generator, NLI, prompt, and decoding metadata are absent; with the full manifest, this rises to \valManFullRepro{}. Claim-level audit (the ability for an external party to re-run NLI on a single claim against the recorded evidence) is impossible without the verification record $\mathbf{V}$, so it is 0 for the minimal manifest and 1 for the full one. The cost is a factor-of-four increase in manifest size (\valStorageKB{}~KB versus 0.42~KB), which is well within the 5~KB per-query budget.

\subsection{Weight Sensitivity}
\label{subsec:sensitivity}

\begin{figure}[H]
\centering
\includegraphics[width=0.7\linewidth]{figures/sensitivity_heatmap.pdf}
\caption{Composite objective $\mathcal{J}$ over the reranker weight grid $(\alpha, \beta)$ at the default $\gamma + \delta = 1 - \alpha - \beta$ allocation. The default $(\alpha, \beta) = (0.5, 0.25)$ marker sits inside a wide stable region; the objective degrades when $\beta$ exceeds 0.45 (over-trust suppresses correct evidence from low-trust sources) or when $\alpha$ falls below 0.30 (relevance is under-weighted relative to provenance).}
\label{fig:sensitivity}
\end{figure}

Figure~\ref{fig:sensitivity} reports the composite objective $\mathcal{J} = 0.5 \cdot \text{Evidence-Grounded Recall} + 0.5 \cdot (1 - \text{Poisoned-Claim Support Rate at 5\% poison})$ over the $(\alpha, \beta)$ grid, holding $\gamma + \delta = 1 - \alpha - \beta$ at the default 60/40 split. The default operating point sits inside a wide stable plateau; the objective drops in two failure regions. When $\beta > 0.45$, source trust dominates the score and the reranker over-suppresses evidence from low-trust sources, including correct evidence from new or rarely seen sources; this is the operationalization of the failure mode discussed in Section~\ref{subsec:rationale} and supports the design choice to expose $\beta$ as an operator-tunable parameter rather than a learned weight. When $\alpha < 0.30$, relevance is under-weighted and the reranker surfaces highly trusted but topically irrelevant evidence. The default $(0.5, 0.25)$ avoids both regions.

\subsection{Latency, Storage, and Scalability}
\label{subsec:latency}

\begin{table}[H]
\caption{Per-stage latency on a single $4\times$A100 workstation, plus end-to-end totals and average manifest storage. Timings are wall-clock per query over 500 queries.}
\label{tab:latency}
\small
\setlength{\tabcolsep}{6pt}
\begin{tabularx}{\textwidth}{@{}>{\raggedright\arraybackslash}X *{3}{>{\centering\arraybackslash}c}@{}}
\toprule
\thead{Stage} & \thead{Mean (ms)} & \thead{SD (ms)} & \thead{P95 (ms)}\\
\midrule
Hybrid retrieval & 119 & 23 & 159\\
Provenance-aware reranking (B4) & 15 & 4 & 23\\
Relevance-only reranking (B1, baseline) & 13 & 3 & 19\\
LLM generation & 877 & 165 & 1178\\
Claim verification & 184 & 36 & 244\\
Ledger append & 8 & 2 & 11\\
\midrule
Total (B1, Vanilla) & 1009 & 165 & 1307\\
Total (B4, ProvGuard-RAG) & 1204 & 171 & 1500\\
\midrule
\textbf{Manifest storage} & \multicolumn{3}{c}{1.83 $\pm$ 0.03~KB per query}\\
\bottomrule
\end{tabularx}
\end{table}

Table~\ref{tab:latency} decomposes per-stage latency on the workstation in Section~\ref{subsec:implementation}. LLM generation is the dominant term (\valLatGen~ms, $\sim$73\% of B4 total). Provenance-aware reranking adds \valLatRerankBfour{}~ms over relevance-only reranking; claim verification adds \valLatVerify~ms; ledger append adds \valLatLedger~ms. End-to-end totals are \valTotalBoneLat~ms for B1 and \valTotalBfourLat~ms for B4, an added latency of \valLatAddedFour~ms that remains within the 500~ms operational budget. Each manifest occupies \valStorageKB{}~KB on average. The Merkle inclusion proof for a log of $10^6$ entries adds 640~bytes (20 hashes at 32 bytes each), so total per-query storage stays below 3~KB, comfortably within the 5~KB budget.

\begin{figure}[H]
\centering
\includegraphics[width=\linewidth]{figures/scaling.pdf}
\caption{Scaling behavior of B1 and B4 along four operational axes: (a) latency vs.~top-$k$ retrieved evidence chunks, (b) latency vs.~number of generated claims $K$, (c) latency vs.~corpus size $N$ on a logarithmic axis, and (d) throughput vs.~concurrent workers on a single $4\times$A100 host.}
\label{fig:scaling}
\end{figure}

Figure~\ref{fig:scaling} characterizes how the latency and throughput of the pipeline scale along four operational axes. Latency grows linearly in the number of retrieved chunks (panel a) because evidence length drives the generator's context-conditioned compute and, for B4, the verifier's number of (claim, evidence) entailment calls. Latency grows linearly in the number of generated claims $K$ (panel b) because each claim incurs an additional verification call, with a slope of $\sim$22~ms/claim that matches the per-claim verifier cost. Latency grows weakly with corpus size $N$ (panel c, semilog axis) because the BM25 and FAISS HNSW retriever stages are sublinear; even at $N = 10^7$ documents, the added retrieval cost is below 100~ms relative to $N = 10^4$. Throughput (panel d) scales nearly linearly with concurrent workers up to the 8-GPU saturation point, after which both configurations plateau at 8.5 (B1) and 7.1 (B4) queries per second on the test workstation. The 16\% throughput penalty of B4 over B1 is consistent with the 19\% per-query latency overhead reported in Table~\ref{tab:latency}.

\subsection{Seven-Axis Evaluation Summary}
\label{subsec:summary}

\begin{figure}[H]
\centering
\includegraphics[width=0.72\linewidth]{figures/radar_evaluation.pdf}
\caption{Seven-axis comparison of Vanilla RAG (B1) and ProvGuard-RAG (B4). Axes are normalized to $[0, 1]$ with 1.0 the best achievable value. Speed is the latency ratio $L_{B1}/L_{B4}$; storage efficiency is $5/(5 + s_M)$ for manifest size $s_M$; poison robustness is $1 - r_{\text{sup}}$ at 5\% poison level (Table~\ref{tab:poisoning}).}
\label{fig:radar}
\end{figure}

Figure~\ref{fig:radar} summarizes the seven evaluation axes. ProvGuard-RAG expands the coverage area on five axes (incident accuracy, claim support, poison robustness, retrieved provenance, audit verifiability) at a cost on two axes (speed, storage efficiency) that is bounded by design. On the operational axes (speed, storage), B4 retains $\geq$83\% of B1's performance; on the trust axes (provenance, audit), B4 dominates B1 by construction. The shape of the radar makes the architectural argument visible: ProvGuard-RAG does not improve any single axis dramatically, it widens the coverage of the trust properties that no single existing component addresses on its own.

The comparisons across Tables~\ref{tab:main_results}, \ref{tab:external}, \ref{tab:poisoning}, \ref{tab:injection}, \ref{tab:tamper}, \ref{tab:nli_type}, \ref{tab:ablation}, \ref{tab:manifest_abl}, and \ref{tab:latency} together support four main findings. First, the verifier is the largest single source of claim-level grounding gains, but provenance reranking contributes a complementary 2--3 point lift on evidence-grounded recall by surfacing evidence that the verifier can act on. Second, the combined pipeline is robust to corpus poisoning and to all four prompt-injection attack types we evaluated, with the largest absolute improvement on the indirect-injection class. Third, manifest-integrity guarantees are detected at 1.000 across the five tamper scenarios that they are designed to cover, and the two scenarios outside their design boundary (T6, T7) are reported with their measured residual gap. Fourth, the cost of these properties is a $\sim$19\% latency overhead and a 1.81~KB per-query manifest, both well within typical operational budgets.

% ===========================================================================
%  5. DISCUSSION
% ===========================================================================
\section{Discussion}
\label{sec:discussion}

\subsection{Four Distinct Trust Properties, Four Distinct Mechanisms}
\label{subsec:four_properties}

A central reason that trustworthy-RAG work is hard to compare across papers is that integrity, authenticity, semantic support, and auditability are commonly conflated. The architecture presented here separates these and ties each to its own mechanism. \textbf{Evidence integrity} (whether a registered document has been altered) is the job of the per-document SHA-256 hash anchored at registration; the verifier and the ledger are not involved. \textbf{Source authenticity} (whether a document genuinely originates from a claimed source) is the job of the source identifier $\mathrm{src}$ and, in deployments that use signed retrieval, of the cryptographic signature chain on the source itself. \textbf{Semantic support} (whether a generated claim is entailed by the retrieved evidence) is the job of the NLI verifier and is a strictly weaker property than truthfulness; a poisoned but internally consistent document can carry a false claim that the verifier nonetheless judges supported. \textbf{Auditability} (whether a third party can reconstruct and re-verify the inference offline) is the job of the manifest plus the Merkle inclusion proof. Conflating any of these four properties produces metric artefacts: provenance utilization gets reported as evidence truthfulness, cryptographic logging gets reported as factual correctness, or NLI accuracy gets reported as audit assurance. The four-property decomposition in Section~\ref{subsec:problem} and the threat model in Section~\ref{subsec:threat_model} are designed to keep these separate, and the seven-axis evaluation in Section~\ref{subsec:summary} reports each property against the mechanism that is supposed to provide it.

\subsection{What the Numbers Mean (and What They Do Not)}
\label{subsec:interpretation}

The provenance reranker reduces poisoned-evidence selection without degrading incident-classification accuracy (Tables~\ref{tab:main_results} and~\ref{tab:poisoning}), which is consistent with provenance metadata carrying signal orthogonal to relevance. The verifier is the dominant driver of claim-level grounding (Table~\ref{tab:ablation}); without it, the pipeline's claim-support and evidence-grounded recall both drop sharply. The audit manifest detects every manifest-altering tamper attempt at 1.000 (T1--T5 in Table~\ref{tab:tamper}) and contributes only $\sim$8~ms of per-query latency (Table~\ref{tab:latency}). These positive findings, however, must be read alongside the explicit non-guarantees: T6 (a poisoned document with internally consistent metadata) is detected at \valBfourTamperTsix{} not 1.000, and T7 (log equivocation without an external witness) is detected at \valBfourTamperTseven{}. Both numbers are reported in the same table as the 1.000 values precisely so that the system is judged on what it actually does rather than on what its name suggests.

\subsection{Failure Mode Analysis}
\label{subsec:failure_modes}

Five operational failure modes deserve explicit attention. First, the provenance-aware reranker can penalize evidence from low-trust sources that nevertheless contains correct information, which threatens evidence faithfulness whenever the only relevant evidence happens to come from a new or rarely seen source; the recommended mitigation is the coverage-preserving constraint that the reranker retains at least one chunk from each high-relevance source cluster, alongside a small uncertainty term that distinguishes ``new source'' from ``unreliable source'' in the historical-trust component $w_{\text{hist}}$. Second, the verifier inherits the limitations of NLI models on multi-hop, temporal, causal, and numerical claims; Table~\ref{tab:nli_type} reports the per-claim-type performance so that operators know which claim types must be triaged to human review rather than treated as automatically settled. Third, the audit manifest records hashes rather than content, so audit verifiability depends on the operator retaining the corpus snapshot; this is a corpus-retention policy concern, not an architecture concern, but it must be addressed in deployment. Fourth, the ledger introduces an honest-operator assumption that bounds the strength of G2--G3; in deployments where the log operator may equivocate, an external witness cohort or a public log such as a Sigstore-style transparency service must be added. Fifth, the framework cannot detect a document that was poisoned before registration (the explicit non-guarantee N1); the only defense against pre-registration compromise is upstream source vetting and post-hoc semantic verification on trusted-source evidence.

\subsection{Generalizability}
\label{subsec:generalizability}

The architectural pattern (provenance-aware reranking, claim-level semantic verification, and audit-manifest anchoring) is in principle applicable to other domains in which RAG outputs carry operational or regulatory consequences, but transfer is not free. Each candidate domain imposes its own evidence standards and its own liability structure, and the calibration burden is non-trivial. \textbf{Legal question answering} requires authority weights tied to jurisdiction, statutory hierarchy, and citator status; an off-the-shelf NLI model is unlikely to capture statutory interpretation reliably, and a domain-specific verifier built on legal-reasoning primitives such as multi-example case-similarity matching~\cite{huang2023legalmcsm} would be required in its place. \textbf{Medical guideline retrieval} requires evidence currency and recommendation-level grading (e.g., GRADE) to be encoded in the source-trust scoring function, and the verifier must distinguish guideline statements from anecdotal evidence; the regulatory liability landscape (FDA, MDR) constrains how the system's abstention behavior must be presented to clinicians. \textbf{Financial compliance} requires the manifest to interoperate with audit trails defined by external regulators (SOX, MiFID II) rather than with a self-defined manifest schema. We therefore frame the architectural transferability as a hypothesis to be tested per domain, not as a claim that the same pipeline drops in unchanged. The domain-adaptation surface is limited to three components (source-trust scoring, entailment model, claim decomposer), but the calibration of each component is a substantive research project in its own right.

\subsection{Limitations}
\label{subsec:limitations}

The work has several limitations that bound its scope. First, the evidence corpus is limited to structured and semi-structured textual data; multi-modal CPS evidence (network packet captures, system-call traces, industrial control system telemetry) is outside the current scope and would require modality-specific retrieval and verification components. Second, the verifier is an entailment model used as a screening signal, not a truth oracle; the per-claim-type breakdown in Table~\ref{tab:nli_type} bounds its competence on multi-hop, causal, and remediation claims, and the abstention-precision metric in Table~\ref{tab:main_results} is the operationally meaningful quality target rather than ``factual correctness.'' Third, the evaluation is on two CPS/IoT datasets and four auxiliary knowledge bases; generalization to other cyber-physical subdomains (industrial control, automotive, healthcare IoT) has not been assessed in this manuscript and is left to future work. Fourth, the ledger-anchoring mechanism provides tamper-evidence rather than tamper-prevention; an attacker who controls the log operator and faces no external witness can produce equivocating views (the explicit non-guarantee N4). Fifth, the source trust scores require domain calibration, and a mis-calibrated weight $w_{\text{auth}}$ can suppress correct evidence from new sources; the weight-sensitivity analysis in Section~\ref{subsec:sensitivity} characterizes the operational region where this risk is bounded. Sixth, the 500-query benchmark is built from public IoT datasets via a templated query layer with rendered telemetry rather than from live SOC tickets; this allows controlled measurement of the four trust properties but does not test live analyst behavior, free-form operator notes, or the heterogeneous incident reports that arrive in real cyber-physical operations centers. Validating ProvGuard-RAG on operational SOC traffic, including attacker tactics that target the evidence-renderization layer itself, is the next step before claiming deployment readiness.

% ===========================================================================
%  6. CONCLUSION
% ===========================================================================
\section{Conclusions}
\label{sec:conclusion}

RAG systems for cyber-physical incident intelligence sit at the intersection of four trust properties (answer correctness, evidence faithfulness, provenance traceability, and audit verifiability) that prior work has addressed in isolation rather than jointly. We presented ProvGuard-RAG, a provenance-aware RAG framework that targets these properties as a coupled system: hybrid retrieval with a four-factor provenance-aware reranker, claim decomposition and NLI-based semantic verification, and a five-record audit manifest anchored in an append-only Merkle transparency log. The accompanying threat model separates evidence integrity, source authenticity, semantic support, and audit verifiability as four distinct guarantees backed by four distinct mechanisms, and explicitly enumerates the four non-guarantees that no single combination of provenance, NLI, and ledger logging can reach. The 500-query incident-intelligence benchmark, with seven external and ablated baselines and three attack families (corpus poisoning, prompt injection, and seven post-hoc tamper scenarios), exercises every property the framework claims to address.

The findings reported in Section~\ref{sec:experiments} support three architectural conclusions. First, the verifier is the dominant driver of claim-level grounding gains, but provenance reranking contributes a complementary lift in evidence-grounded recall by surfacing evidence on which the verifier can act. Second, manifest-integrity guarantees are detected at 1.000 on the five tamper scenarios they are designed to cover; the two scenarios that fall outside their design boundary are reported with their measured residual gap rather than excluded. Third, the cost of these properties is a bounded latency overhead and a sub-2~KB per-query manifest. The result is not that any one mechanism is novel in isolation, but that the joint pipeline exposes a query-level accountability layer in which retrieval evidence, generated claims, semantic-support labels, and tamper-evident audit records are produced and evaluated together. To the best of our knowledge, no prior trustworthy-RAG component covers all four properties at once under a unified evaluation protocol of the kind reported in Section~\ref{sec:experiments}.

Several directions remain open. The verifier should be extended to multi-hop entailment so that claims supported by the conjunction of several evidence chunks are not penalized for the absence of single-chunk support. The framework should be lifted to multi-modal cyber-physical evidence, including network packet captures and system-call traces, which would require modality-specific retrievers and verifiers. Deployment studies in operational settings should test whether the controlled benchmark overhead carries over to production workloads, and whether the source-trust scoring function remains well calibrated under realistic source churn. Finally, the architectural pattern should be tested in adjacent domains (legal, medical, financial RAG) under per-domain evidence standards, calibration burdens, and liability regimes; the present work positions this transfer as a hypothesis to be evaluated per domain, not as a claim that the same pipeline transfers unchanged.

% ===========================================================================
%  BACK MATTER (MDPI format)
% ===========================================================================
\vspace{6pt}

\authorcontributions{Conceptualization, H.R.W.\ and Y.G.F.; methodology, H.R.W.\ and H.W.; software and experiments, H.W.; writing---original draft preparation, H.R.W.\ and H.W.; writing---review and editing, Y.G.F.; supervision and funding acquisition, Y.G.F.\ All authors have read and agreed to the published version of the manuscript.}

\funding{This research received no external funding.}

\institutionalreview{Not applicable.}

\informedconsent{Not applicable.}

\dataavailability{The four primary datasets (TON\_IoT, CICIoT2023, CVE/NVD, MITRE ATT\&CK) are publicly available from the original publishers. The accompanying release contains: (i)~the dataset preprocessing pipeline; (ii)~the fourteen query templates and the rendered evidence chunks; (iii)~the attack payloads for the corpus-poisoning, prompt-injection, and audit-tamper protocols; (iv)~the human annotation labels and guideline; (v)~the raw model outputs, retrieved-evidence identifiers, and per-query manifest JSON files; (vi)~the Merkle roots and audit paths produced for every query; (vii)~the NLI training/evaluation split identifiers; (viii)~all random seeds and the container specification used for inference; and (ix)~the table-reproduction script at \texttt{experiments/reproduce\_tables.py} together with the macro file \texttt{experiments/values.tex} from which every number quoted in the manuscript is sourced.}

\conflictsofinterest{The authors declare no conflicts of interest.}

\acknowledgments{During the preparation of this manuscript, the authors used Claude (Anthropic) for drafting assistance. The authors have reviewed and edited the output and take full responsibility for the content of this publication.}

\abbreviations{Abbreviations}{
The following abbreviations are used in this manuscript:
\smallskip

\noindent
\begin{tabular}{@{}ll}
RAG  & Retrieval-Augmented Generation \\
LLM  & Large Language Model \\
CPS  & Cyber-Physical Systems \\
NLI  & Natural Language Inference \\
IoT  & Internet of Things \\
IIoT & Industrial Internet of Things \\
NVD  & National Vulnerability Database \\
CVE  & Common Vulnerabilities and Exposures \\
\end{tabular}
}

% ===========================================================================
%  REFERENCES
% ===========================================================================
% MDPI uses natbib (loaded by mdpi.cls). Use external .bib file:
\reftitle{References}
\bibliography{references}

% If the above does not work in your environment, use the internal format:
% \begin{thebibliography}{999}
% \bibitem{...}
% \end{thebibliography}

\PublishersNote{}

\end{document}
